\documentclass[11pt,a4paper]{article}
\usepackage[T1]{fontenc}
\usepackage[utf8]{inputenc}
\usepackage{lmodern}
\usepackage{amsmath,amssymb,amsthm,mathtools}
\usepackage[margin=27mm,headheight=14pt]{geometry}
\usepackage{microtype,booktabs,array,longtable,enumitem}
\usepackage{xcolor}
\usepackage{fancyhdr}
\usepackage{hyperref}
\usepackage[nameinlink,noabbrev]{cleveref}
\hypersetup{colorlinks=true,linkcolor=blue!45!black,citecolor=blue!45!black,
 urlcolor=blue!45!black,pdftitle={Sharp Matrix-Scaling Stability over Surreal and Surcomplex Fields},
 pdfauthor={AI-assisted research manuscript prepared for Vladimir Reshetnikov}}
\pagestyle{fancy}
\fancyhf{}
\fancyhead[L]{\small Sharp matrix-scaling stability}
\fancyhead[R]{\small Surreal and surcomplex fields}
\fancyfoot[C]{\thepage}
\renewcommand{\headrulewidth}{0.3pt}
\setlength{\parskip}{3pt}
\setlength{\parindent}{1.2em}
\setlength{\emergencystretch}{2em}
\newcolumntype{L}[1]{>{\raggedright\arraybackslash}p{#1}}
\setlist{itemsep=3pt,topsep=5pt}
\numberwithin{equation}{section}
\newtheorem{theorem}{Theorem}[section]
\newtheorem{proposition}[theorem]{Proposition}
\newtheorem{lemma}[theorem]{Lemma}
\newtheorem{corollary}[theorem]{Corollary}
\theoremstyle{definition}
\newtheorem{definition}[theorem]{Definition}
\newtheorem{example}[theorem]{Example}
\newtheorem{remark}[theorem]{Remark}
% Added during assembly: statements share the theorem counter, and without these
% aliases cleveref names every lemma, corollary, definition, example and
% remark a "theorem".
\AddToHook{env/proposition/begin}{\crefalias{theorem}{proposition}}
\AddToHook{env/lemma/begin}{\crefalias{theorem}{lemma}}
\AddToHook{env/corollary/begin}{\crefalias{theorem}{corollary}}
\AddToHook{env/definition/begin}{\crefalias{theorem}{definition}}
\AddToHook{env/example/begin}{\crefalias{theorem}{example}}
\AddToHook{env/remark/begin}{\crefalias{theorem}{remark}}
\newcommand{\R}{\mathbb R}
\newcommand{\C}{\mathbb C}
\newcommand{\Q}{\mathbb Q}
\newcommand{\Z}{\mathbb Z}
\newcommand{\N}{\mathbb N}
\newcommand{\No}{\mathbf{No}}
\newcommand{\OO}{\mathcal O}
\newcommand{\mm}{\mathfrak m}
\newcommand{\GG}{\mathcal G}
\newcommand{\TT}{\mathcal T}
\newcommand{\NN}{\mathcal N}
\newcommand{\supp}{\operatorname{supp}}
\newcommand{\im}{\operatorname{im}}
\newcommand{\diag}{\operatorname{diag}}
\newcommand{\lc}{\operatorname{lc}}
\newcommand{\rank}{\operatorname{rank}}
\newcommand{\res}{\operatorname{res}}
\newcommand{\od}{\mathbin{\odot}}
\newcommand{\eps}{\varepsilon}
\newcommand{\ph}{\varphi}
\newcommand{\pin}{048b72cf7cbfc8ab246e4f73788c10460cb3f6e0}
\newcommand{\shortpin}{048b72c}
\newcommand{\repo}{\url{https://github.com/VladimirReshetnikov/Surreal}}
\newcommand{\valvec}[1]{\min_{e\in E}v(#1_e)}

\begin{document}
\hypersetup{pageanchor=false}
\begin{titlepage}
\vspace*{7mm}
\noindent{\small RESEARCH MANUSCRIPT \hfill 22 September 2026}
\vspace{11mm}
\begin{center}
{\LARGE\bfseries Sharp Matrix-Scaling Stability\\[4pt]
over Surreal and Surcomplex Fields}\par
\vspace{6mm}
{\large Spanning-tree gaps, infinitesimal normalization,\\[3pt]
and exact multiscale propagation}\par
\vspace{8mm}
{\normalsize Prepared for Vladimir Reshetnikov}
\end{center}
\vspace{7mm}
\begin{abstract}
We study diagonal matrix scaling in set-sized Hahn fields
$F=\R((t^\Gamma))$ and $K=\C((t^\Gamma))$, with no rank or divisibility
assumption on the ordered value group for the local theory. A positive
matrix may have entries at arbitrarily different infinitesimal scales.
Nevertheless, normalization to its fixed row and column sums admits a
unique support-controlled surcomplex analytic branch on the entire
relative infinitesimal polydisc.

The main quantitative invariant is the \emph{spanning-tree deletion gap}.
For an edge $e$ of the bipartite support graph, let $\tau$ be the minimum
sum of entry valuations over spanning trees and let $\tau_e^-$ be the same
minimum over trees avoiding $e$. Then $\kappa_e=\tau_e^--\tau$ is the
uniform gain in relative valuation at entry $e$, and is optimal when
$\Gamma\ne\{0\}$: an input change
of valuation at least $\delta>0$ produces an output change of valuation
at least $\delta+\kappa_e$. The assertion is a nonlinear two-point
estimate, not merely an infinitesimal derivative formula. Every change
in just entry $e$ attains equality at that entry; a bridge is fixed
exactly. All positive-degree Taylor coefficients satisfy the same
entrywise valuation bounds.

A tree interpolation identity and a secant reweighting of the exponential
prove the sharp estimate. A separate finite-generator support argument
licenses evaluation at arbitrary rank, where integrality of formal
coefficients alone is insufficient. For a tridiagonal bistochastic
chain we obtain an exact nonlinear propagation law: valuation gains add
along consecutive links. We also give a real-linear-constraint
extension, an unramified local descent statement, explicit examples,
and exact finite verification programs.
\end{abstract}
\vfill
\noindent\small\textbf{Status and priority.}
The main theorems are proved in this manuscript. Their originality is a
\emph{candidate} conclusion of a targeted, non-exhaustive comparison with
published work and the repository at commit \texttt{\shortpin}.
No named long-standing open problem is claimed solved, and no refereeing
or complete machine verification is claimed. Known scaling, projection,
and spanning-tree results are explicitly credited.
This is an AI-assisted research draft. Its 16 standard statements are
inventoried in the collection's ledger, but none has an Implementation mapping;
no Lean proof coverage is claimed here. It is built from a single
manuscript; \cref{scale:sec:provenance} records how that manuscript was
placed in the collection and what was changed.
\end{titlepage}
\hypersetup{pageanchor=true}
\pagenumbering{roman}

\tableofcontents
\clearpage
\pagenumbering{arabic}

\section{The problem and the contribution}\label{scale:sec:intro}

\subsection{Why this is more than scalar extension}
Let $P=(p_{ij})$ be a finite matrix with positive surreal entries on a
fixed support. Change its nonzero entries by relative infinitesimals,
and then multiply rows and columns by positive factors so that the
original margins are restored. How accurately are the individual
entries of the normalized matrix determined?

A direct inverse-Jacobian bound can suggest a loss of precision when
some entries are extremely small. A different phenomenon occurs for
fixed margins. The normalized entries never lose relative valuation,
and certain entries gain a precisely computable additional amount.
The answer depends on which edges are indispensable to
minimum-valuation spanning trees, rather than on a smallest eigenvalue
used as a scalar condition number.

This is a genuinely multiscale question. For example, in
$\Gamma=\Z^2$ with lexicographic order, put
$\alpha=(0,1)$ and $\beta=(1,0)$. Then $n\alpha<\beta$ for every
ordinary positive integer $n$. An argument requiring $n\alpha$ to
become arbitrarily large cannot justify a Hahn expansion at
$t^\alpha$. The proofs below retain the entire ordered group and use
strong summability, not an implicit replacement by a real-valued norm.

The word \emph{surcomplex} means the complexification $\No[i]$.
For rigorous constructions we first work in the set-sized field
$K=\C((t^\Gamma))$. The final passage to surreal and surcomplex numbers
uses the normal-form embedding and introduces no class-sized sums;
see \cref{scale:sec:surreal}. The local existence theorem is actually
slightly stronger than a real-positive statement: each base weight may
be surcomplex, provided its leading coefficient is positive real.

\subsection{An overview of the exact result}
For a connected bipartite graph $\GG=(R\sqcup C,E)$ and edge weights
$p_e$, write $c_e=v(p_e)$. Let $\TT$ be its finite set of spanning trees.
Define
\begin{equation}\label{scale:eq:overviewgap}
 \tau=\min_{T\in\TT}\sum_{a\in T}c_a,\qquad
 \tau_e^-=\min_{\substack{T\in\TT\\e\notin T}}\sum_{a\in T}c_a,
 \qquad \kappa_e=\tau_e^--\tau.
\end{equation}
An empty minimum is $+\infty$. Thus a bridge has $\kappa_e=+\infty$.
For $f\in\mm_K^E$ let the perturbed kernel be
$A_e(f)=p_e\exp(f_e)$, using only the infinitesimal exponential.
Let $q(f)$ denote its normalization to the margins of $p$ on the
relative infinitesimal branch. The central estimate is
\begin{equation}\label{scale:eq:overviewbound}
 v\!\left(\frac{q_e(f)}{q_e(g)}-1\right)
 \geq \kappa_e+\min_{a\in E}v(f_a-g_a),\qquad f,g\in\mm_K^E.
\end{equation}
When $f-g=s\mathbf e_e$, where $0\ne s\in\mm_K$, equality holds in
\eqref{scale:eq:overviewbound} for a nonbridge $e$. If $\Gamma\ne\{0\}$,
such an $s$ exists (take $s=t^\gamma$ with $\gamma>0$), so no larger
uniform gain can replace $\kappa_e$. When $\Gamma=\{0\}$, the
infinitesimal domain is the singleton $\{0\}$; the coefficientwise
sharpness below still holds, but evaluated optimality has no witness.
The equality for nonzero changes holds even when
minimum trees are tied and when the perturbations have nonreal
coefficients.

An equivalent combinatorial formula needs no tree enumeration. If
$e$ joins $u$ to $v$, then
\begin{equation}\label{scale:eq:overviewpath}
 \kappa_e=\max\left\{0,
 \min_{\pi:u\leadsto v\text{ in }\GG\setminus e}
       \max_{a\in\pi}c_a-c_e\right\}.
\end{equation}
The value is $+\infty$ if no alternative path exists. The minimax
path formula is familiar minimum-spanning-tree sensitivity. Its role
as the \emph{optimal nonlinear normalization gain}, including all
Taylor orders, is the contribution investigated here.

\subsection{What is proposed as new, and what is not}
\begin{center}
\begin{tabular}{@{}L{0.27\textwidth}L{0.65\textwidth}@{}}
\toprule
Result & Status in this manuscript\\
\midrule
Scaling existence over real closed fields &
Classical real scaling plus first-order transfer; included as background.\\[3pt]
Weighted constraint projection and tree formulas &
Known linear-algebraic and transfer-current ingredients; proved directly
for transparency.\\[3pt]
Whole-infinitesimal-polydisc normalization &
Candidate original arbitrary-rank, positive-leading surcomplex theorem,
with an explicit strong-summability proof.\\[3pt]
Sharp tree-deletion gain &
Candidate original nonlinear two-point law, coefficient bounds, and exact
one-coordinate sharpness.\\[3pt]
Propagation through a chain &
Candidate original exact nonlinear valuation formulas at all links and
diagonal entries.\\[3pt]
Real-linear-constraint extension &
A proved extension replacing trees by bases; no new matroid combinatorics
is claimed.\\
\bottomrule
\end{tabular}
\end{center}
The first derivative alone is \emph{not} a novelty claim. Nor is the
mere existence of generalized series for a one-parameter scaling
problem. Both have substantial precedents, described next.

\subsection{What this report does not claim}\label{scale:sec:nonclaims}
Each limitation below is also stated where it applies. They are collected
here so that none of them is lost when the results are quoted elsewhere.
\begin{enumerate}[label=(N\arabic*),leftmargin=3.2em]
\item \emph{Priority and status.} The whole-polydisc normalization, the
sharp tree- and basis-gap laws and the chain propagation formulas are
\emph{candidate} original results after a targeted, non-exhaustive
comparison with published work and with the repository at its pin
(\cref{scale:sec:comparison}). Publication priority, absence from every
repository file and independent expert validation are not certified. The
repository audit was not a line-by-line review of every article or Lean
declaration, and its indexed searches were unusable as negative evidence.
No named long-standing open problem is claimed solved. The report is an
AI-assisted draft; it has not been refereed; none of its results has a Lean
Implementation mapping in the ledger, and \cref{scale:sec:verification} does not assert that the
Lean modules it outlines exist.
\item \emph{Credited, not claimed.} Existence, uniqueness and ordinary
real continuity of prescribed-margin scaling (Idel's review and classical
scaling theory); generalized-series expansions of a one-parameter scaling
problem, so that no first use of Hahn-type series in matrix scaling is
claimed (Sharify--Gaubert--Grigori); the first-order implicit
differentiation formula (Eisenberger et al.); the determinant formula, its
tree interpretation and transfer-current representations (Burton--Pemantle);
the minimax path formula~\eqref{scale:eq:overviewpath} and the
minimum-spanning-tree exchange arguments of \cref{scale:sec:bottleneck};
Hahn and surreal normal-form background (Gonshor, Berarducci--Mantova);
Higman's word theorem; real-closed transfer (Tarski). The real-linear-constraint
extension (\cref{scale:thm:linear}) claims no new matroid combinatorics.
\item \emph{No convergence assertion.} Neither convergence of the alternating
Sinkhorn iteration in an arbitrary-rank valuation topology, nor its iteration
complexity, nor convergence of set-indexed sequences in the full surreal fine
topology is claimed. Strong sums are formed coefficientwise; strong
summability alone does not imply convergence of partial sums
(\cref{scale:sec:boundaries}). The remainder certificate
(\cref{scale:cor:remainder}) does not promise that a finite Taylor degree
reaches every requested precision at higher rank.
\item \emph{What ``analytic'' means.} The normalization branch is given by a
strongly evaluated formal series on the relative infinitesimal polydisc. No
classical holomorphic manifold structure in the full surreal fine topology
is asserted. No exponential or logarithm at an infinite surcomplex argument,
no global surcomplex exponential and no surreal derivation is used; Taylor
coefficients are formal derivatives with respect to the edge parameters.
\item \emph{The hypotheses matter.} Positive leading coefficients of the
base weights are sufficient, not claimed to be necessary; without them
integrality and uniqueness can fail
(\cref{scale:ex:phasejacobian,scale:ex:phasebranches}). The field $K$ is not
ordered, and ``positive-leading'' refers only to the leading coefficient.
The margins are fixed; if they move, even by relative infinitesimals, the
bridge rigidity fails (\cref{scale:sec:boundaries}). The support is fixed.
\item \emph{What ``sharp'' means.} For $\Gamma\ne\{0\}$,
$\kappa_e$ is the largest gain valid uniformly over infinitesimal inputs.
Formal coefficient sharpness also holds for the trivial group;
equality is not asserted for every input direction
(\cref{scale:thm:sharp} and the remark after it). Optimality of the exponent
$N$ in \cref{scale:prop:cutbound} is not claimed.
\item \emph{Which costs.} The tree costs of the sharp theorem are valuations
of the normalized base matrix. No equally sharp formula from the valuations
of an unscaled kernel is claimed (\cref{scale:cor:positivebranch}).
\item \emph{Divisibility and real closedness.} The local theory assumes no
divisibility of $\Gamma$. The global positive interpretation
(\cref{scale:prop:global,scale:cor:positivebranch} and the global paragraph
of \cref{scale:sec:surreal}) uses real closedness, reached through the
divisible hull and first-order transfer; the real optimization proof is not
repeated inside a non-Archimedean field, and no non-Archimedean compactness
or global logarithm is used. The descent statement \cref{scale:cor:descent}
is local: global scaling over a nondivisible group can require ramification
\eqref{scale:eq:ramificationexample}.
\item \emph{Real constant constraints.} \Cref{scale:thm:linear} needs a real
constant constraint matrix. The same basis-cost formula is not asserted for a
Hahn-valued constraint matrix.
\item \emph{Surreal transport.} The surreal and surcomplex statements are
obtained by set-sized support-group localization. No class-sized sum, net or
compactness argument is used, and no proof of a local result uses spherical
completeness or Banach contraction.
\item \emph{Computation.} The workflow of \cref{scale:sec:verification} is an
exact-operation specification, not a universal computable representation of
surreal numbers: arbitrary Hahn coefficients need not come with an effective
leading exponent, equality test or support oracle. The amplitude $C_e$ needs
all minimizing trees when there are ties. The finite checks of
\cref{scale:sec:verification} do not prove strong summability, arbitrary-rank
existence or surreal transport, and a successful compilation is not a proof
check.
\item \emph{Open questions.} The three questions of \cref{scale:sec:verification}
(a sharp two-edge propagation formula on a general support graph, of which the
chain theorem solves one family; a classification of the surcomplex base
weights with a unique integral whole-polydisc branch; and the convergence and
complexity of alternating normalization over higher-rank valued fields,
together with a global surcomplex scaling theorem for arbitrary complex
kernels) are not settled.
\end{enumerate}
The delivered files \texttt{SOURCE\_AUDIT.md} and \texttt{PROOF\_AUDIT.md}
state the audit boundaries in more detail.

\subsection{Relation to other reports of the collection}\label{scale:sec:neighbours}
The proofs in this report are given directly; the neighbouring reports below
supply comparisons. These were added during editorial assembly at
\texttt{ed88b8f} and checked again against the source labels at
\texttt{034ab96} during the main-text review. Historical search claims
remain tied to their snapshots in \cref{scale:sec:comparison}.

\paragraph{Markov generators at every scale.}
The Markov report \cite{RepoMarkov}, held in
\nolinkurl{docs/surreal/markov-generators-at-every-scale/}, and this one
share one mechanism: a positive combinatorial
expansion whose valuation is a minimum over trees or forests and whose
leading coefficient is a positive sum over the minimizers. There,
Proposition~3.1 (\nolinkurl{markov:prop:forest}, the directed matrix-forest
theorem) writes $\det(sI+L)$ and $\operatorname{adj}(sI+L)$ as sums of
positive in-forest weights, and equation~(3.10) there
(\nolinkurl{markov:eq:profile-properties}) reads off the minimal forest weight
and its amplitude. Its Theorem~4.1 (\nolinkurl{markov:thm:leading}) gives the
valuation of a resolvent entry as a minimum over forests with a prescribed
root, minus the unrestricted minimum. Here \cref{scale:thm:tree} writes
$\det L_p$ as the positive tree sum $Z(p)$, \eqref{scale:eq:Hdiagonal} writes
a diagonal projection entry as $Z_e^-(p)/Z(p)$, and the gap
$\kappa_e=\tau_e^--\tau$ is again a restricted minimum (trees avoiding $e$)
minus the unrestricted one, with the quotient of amplitudes
$C_e$ in \eqref{scale:eq:leadinggain}. As polynomial identities in the
weights, symmetric rates $q_{uv}=q_{vu}=p_e$ on the support graph make the
in-forests of that report with $N-1$ edges and a fixed root exactly the
spanning trees oriented towards that root, so its top forest coefficient
$\sigma_{N-1}$ equals $N\,Z(p)$. This report does not use that identity.

The differences are as follows. (i) \emph{Objects.} There, the row Laplacian
of a directed graph, in general neither symmetric nor normal, and the
normalized resolvent, a rational function of the rates and of a scale
parameter, with forests of every size. Here, the symmetric reduced Laplacian
$BW_pB^{\mathsf T}$ of a bipartite support graph, only spanning trees, no
scale parameter, and a \emph{nonlinear} fixed-margin normalization obtained
from a support-controlled power series (\cref{scale:thm:existence}).
(ii) \emph{Stability.} Theorem~8.1 there (\nolinkurl{markov:thm:stability}) turns
relative rate errors of valuation at least $\delta$ into relative
resolvent-entry errors of valuation at least $\delta$, uniformly in the
resolvent parameter, and Example~8.3 there shows that no gain is possible in
general. \Cref{scale:thm:sharp} turns relative kernel changes of valuation at
least $\delta$ into normalized-entry changes of valuation at least
$\delta+\kappa_e$, with equality for one-coordinate changes; the gain comes
from the fixed margins and fails when the margins move
(\cref{scale:sec:boundaries}). The two theorems concern different maps, and
neither implies the other. Both need a fixed support: a new edge is not a
relative perturbation (Remark~8.4 there). (iii) \emph{Hypotheses.} That
report assumes a divisible value group throughout and real positive rates.
The local theory here assumes no divisibility and allows surcomplex
positive-leading base weights with complex perturbations. (iv) The
crossover hierarchy, effective generators, rank-one realization and
eigenvalue amplitudes of that report have no counterpart here. The whole-polydisc
branch, the deletion gap and its bottleneck formula, and the chain
propagation law of this report have no counterpart there.

\paragraph{Surcomplex spectral theory.}
The spectral report \nolinkurl{docs/surcomplex/spectral-theory/}
\cite{RepoSpectral} shares three ingredients with this one.
\begin{enumerate}[label=(\alph*)]
\item \emph{Positive Cauchy--Binet sums.} Its Theorem~8.1
(\nolinkurl{thm:gram}, a label without report prefix) writes the coefficients of
$\det(I+sA^*A)$ as sums of squared moduli of minors, so that no leading
cancellation occurs. \Cref{scale:thm:tree} expands
$\det(BW_pB^{\mathsf T})$ by Cauchy--Binet as
$\sum_T(\det B_T)^2w_T(p)$, and \cref{scale:thm:linear} has the same form
with real factors $(\det B_T)^2>0$. When the weights are real and positive
and have square roots in the field, $BW_pB^{\mathsf T}=A^*A$ with
$A=W_p^{1/2}B^{\mathsf T}$, and $\det L_p$ is the top Gram coefficient of that
theorem. The proof here needs neither square roots, which a nondivisible
value group need not supply, nor real weights: it applies Cauchy--Binet to
the product $B\cdot W_p\cdot B^{\mathsf T}$ with positive-leading surcomplex
weights. The interpolation identity~\eqref{scale:eq:treeaverage} for the
projection itself has no counterpart there.
\item \emph{The support engine.} \Cref{scale:lem:neumann} and Lemma~9.4
there (\nolinkurl{spec:lem:neumann}) state the same positive-support word
finiteness, both from Higman's theorem. \Cref{scale:lem:evaluation} with no
generators ($R_0=\C$) gives the strong-evaluation part of Lemma~9.5 there
(\nolinkurl{spec:lem:evaluate}); \cref{scale:lem:evaluation} allows coefficients
in a finitely generated subalgebra of $\OO$, and \cref{scale:ex:badformal}
shows why arbitrary coefficients in $\OO$ are not allowed. Example~9.2 there
(\nolinkurl{spec:ex:nonconvergence}) and the last subsection of
\cref{scale:sec:boundaries} give the same failure of partial-sum convergence,
with the elements $(0,1)$ and $(1,0)$ of $\Z^2_{\mathrm{lex}}$.
Strong summability alone does not imply such convergence. Each report keeps its own proofs.
\item \emph{No new exponents.} Its Theorem~9.18 (\nolinkurl{spec:thm:svd}) keeps
singular value decompositions in the subgroup generated by the input
supports, as \cref{scale:cor:descent} keeps the local normalization in
$\C((t^\Lambda))$. Its Section~10 computes when positive matrix roots need
new exponents. The global ramification example
\eqref{scale:eq:ramificationexample} here is computed directly and is not
derived from that section.
\end{enumerate}
What is not shared: that report is about eigenvalues, singular values and
Hermitian or normal matrices, and its perturbation results (for instance
Theorem~5.2 there, \nolinkurl{thm:weyl}) bound changes of eigenvalues and
singular values. No eigenvalue or singular value appears in this report. The
gain $\kappa_e$ is a combinatorial invariant of the support graph and the
entry valuations, not a smallest singular value used as a condition number
(\cref{scale:sec:intro}), and singular-value stability supplies neither the
whole-polydisc support argument nor the tree-deletion sharpness
(\cref{scale:sec:comparison}).

\paragraph{Prony reconstruction.}
The nonuniform certificate of the Prony reconstruction report
(Theorem~8.1 there, \nolinkurl{prony:thm:graph}) is also a minimum-cost path
computation in an ordered group. Its graph is a directed graph of moment
interactions. Its sufficient reconstruction certificate requires positive
cycle weights and the strict-ball inequalities
$\rho_i-E_i>\delta_i$ for every $i$; its path minima then bound node errors.
Those letters are local to that report. It involves no spanning trees and no normalization, and its moment
precisions $\kappa_k$ are unrelated to the gaps $\kappa_e$.

\subsection{Words and symbols used differently elsewhere in the collection}
\label{scale:sec:words}
Several words and letters of this report mean other things in other reports.
Each has one meaning here.
\begin{itemize}
\item \emph{Scaling} and \emph{normalization} mean diagonal row and column
scaling of a matrix to prescribed margins, and its local branch
$\NN_p$, $q_p$. They are not the monomial rescalings or polynomial
normalizations of other reports. \emph{Scale}, as in the title, means a
valuation.
\item \emph{Positive}, for a matrix, means entries positive on the support
in the ordered field $F$; \emph{positive-leading}
(\cref{scale:def:positiveleading}) is the only positivity notion in $K$,
which is not ordered. Neither means positive definite, the sense of the
spectral report.
\item \emph{Gap} means the spanning-tree deletion gap $\kappa_e$. It is not a
spectral gap and not a surreal gap (a Dedekind section of $\No$) in the sense
of the genetic-gap report. The letter $\kappa$ also denotes moment
precisions in the Prony report and the valuation gain $v(q)$ of an inverse
branch in the expanding-polynomial-dynamics report; those are unrelated.
\item \emph{Tree} means an undirected spanning tree of the support graph. It
is not an in-forest of the Markov report, nor a cluster, separation or
splitting tree of the polynomial, Prony or spectral reports.
\item \emph{Branch} means the unique relative-infinitesimal solution
$\NN_p$ on $\mm^E$, and \emph{analytic} means given there by a strongly
evaluated formal series with coefficients in $R_p$. These are not the
germ-analytic or coherent-section classes of the analysis report, and the
branch is not a branch selector in the computational sense of
\nolinkurl{docs/NOTATION.md}.
\item \emph{Gauge factors} are the diagonal row and column factors of
\cref{scale:cor:gauge}; neither the noise gauge of the Prony report nor a
differential gauge equivalence.
\item \emph{Bistochastic} means that all row and column sums are $1$.
\item $p$ denotes base weights and $q_p(f)$ normalized weights; $P$ and $Q$
are the corresponding matrices. In the Markov report $q_e$ are rates and $P_i$
idempotents. $L_p=BW_pB^{\mathsf T}$ is a symmetric reduced weighted Laplacian,
not that report's row Laplacian. $H_p=I-\Pi_p$ is a projection, not a
Hermitian matrix $H$ as in the spectral report, and $\Pi_p$ is unrelated to
the exponent matrix $\Pi$ of the theta-series material in the Hahn--Tate
report. $s$ is a perturbation scalar, not a resolvent parameter.
\item Local letters are reused with local meanings, each defined where it is
used: $\delta$ (the least valuation of an input change in
\cref{scale:sec:intro,scale:sec:gain,scale:sec:chain}, the group element
$(0,1)$ in \cref{scale:ex:badformal} and \cref{scale:sec:boundaries}, and a
positive parameter in the examples of \cref{scale:sec:boundaries});
$C$ (the column set, the Cramer matrices $C_k$, the amplitude $C_e$, the
cycle basis of \cref{scale:sec:chain}); $R$ (the row set, $R(z)=\exp z-1-z$,
the rings $R_0$ and $R_p$); $D$ (the diagonal factors $D_r,D_c$, the
continuants $D_{a,b}$, a truncation degree). The formal series of
\cref{scale:lem:evaluation,scale:ex:badformal} is written $\Phi(X)$, so that
$F$ always denotes the real Hahn field ($F_0$ in \cref{scale:prop:global} is
an arbitrary real closed field).
\item The valuation is $v(x)=\min\supp x$, and the surreal embedding is
$t^\gamma\mapsto\omega^{-\gamma}$ \eqref{scale:eq:surrealembedding}, as
elsewhere in the collection. \emph{Surcomplex} means $\No[i]$.
\end{itemize}

\subsection{Provenance of this report}\label{scale:sec:provenance}
This report is built from one manuscript, \emph{Sharp Matrix-Scaling
Stability over Surreal and Surcomplex Fields}, dated 22 September 2026. It
arrived as the archive \nolinkurl{surreal_matrix_scaling.zip}, one of nine
archives committed on arrival at \texttt{330aa79}. The archive held the LaTeX
source \texttt{surreal\_matrix\_scaling.tex}, its PDF, a README, the audit
files \texttt{SOURCE\_AUDIT.md} and \texttt{PROOF\_AUDIT.md}, the program
\texttt{code/verify.py} with its recorded output and console record, a build
record with its console log, a requirements file and a build script. It was
placed as raw source at commit \texttt{30dfb4f}; editorial assembly at
\texttt{ed88b8f} added the report PDF and the material described below. Its
own repository comparison is pinned to \texttt{\pin}
(\cref{scale:sec:comparison}).

There is no second source, so there was no merge and no choice between
sources. Nothing was selected out of a larger body of work. Every theorem,
proof, example and disclaimer of the manuscript is kept. The report is a
member of the surreal family because its positive matrices, the global
scaling theory and the tree costs are real Hahn objects; the surcomplex
branch of \cref{scale:thm:existence} extends that real theory.

Raw placement and subsequent editorial assembly made these changes.
\begin{enumerate}[label=(\alph*)]
\item On raw placement, the source was renamed \texttt{article.tex}. The two audit files, the
program and everything in \texttt{data/} are kept byte-identical. The build
script was moved to \texttt{code/} and the requirements file to
\texttt{data/}. The delivered PDF and README are not shipped.
\item All 88 labels of the manuscript already carried the prefix
\texttt{scale:} and are unchanged. The labels added during assembly carry
the same prefix.
\item \Cref{scale:sec:nonclaims,scale:sec:neighbours,scale:sec:words,scale:sec:provenance}
were added during assembly, and so were \cref{scale:rem:support,scale:rem:trees}, which
point to other reports. They were placed so that every theorem, equation
and section number of the manuscript is unchanged. The title-page status
paragraph now states the draft status.
\item Repository statements that had gone stale were corrected, with the pin
kept as provenance (the paragraph ``Since the pin'' in
\cref{scale:sec:comparison}, and \cref{scale:app:build}, whose build
instructions now describe this directory). \Cref{scale:sec:verification}
records a rerun of the checks during assembly.
\item The formal series written $F(X)$ in \cref{scale:lem:evaluation} and
\cref{scale:ex:badformal} was renamed $\Phi(X)$, since $F$ is the real Hahn
field.
\item All numbered statements share one counter. In the delivered build,
every cross-reference to a proposition, lemma, corollary, definition,
example or remark printed the word ``theorem''. Six lines in the preamble now make each
reference print the right name; no number changed.
\end{enumerate}
These assembly changes left the mathematics unchanged. The subsequent
main-text review of Sections 1--14 expanded proofs and corrected the scope
and comparisons: it distinguished evaluated sharpness from formal
coefficient sharpness at the trivial group, specified strict feasibility
on the whole support, corrected two published theorem citations, and
separated strong summability from partial-sum convergence. The delivered
audit, code and data remain unchanged. The review is mathematical reading,
not Lean coverage or a literature-priority certificate; the collection's
\nolinkurl{docs/REVIEW.md} records its scope.

\section{Literature and repository comparison}\label{scale:sec:comparison}

\subsection{The closest published precedents}
The prescribed-margin matrix-scaling problem and its support conditions
are classical. Idel's review \cite{Idel} is a useful entry point;
its Theorem~3.1 gives a general scaling criterion, and Theorem~4.5
(p.~26 of the cited v1) gives continuity on the fixed prescribed-margin
pattern class. The present article does not reclaim
existence, uniqueness, or ordinary real continuity.

Sharify, Gaubert and Grigori \cite{SGG}, Theorem~2.4 (p.~6 of v2),
prove that for a nonnegative square matrix $A$ with total support,
the bistochastic scaling of its Hadamard deformation $A^{(p)}$ has
absolutely convergent generalized Dirichlet expansions in $t=\exp(-p)$.
Their Theorem~2.3 gives convergence to the unique maximum-entropy point
of the optimal-assignment face.
This is an important prior connection between matrix scaling,
generalized series, and tropical methods. Our question is different:
we fix an arbitrary-rank Hahn matrix, allow arbitrary strongly defined
relative infinitesimal changes, and identify the exact edgewise
nonlinear valuation gain. We make no claim to the first use of Hahn-type
series in matrix scaling.

Eisenberger, Toker, Leal-Taix\'e, Bernard and Cremers \cite{Eisenberger}
derive an implicit differentiation framework for Sinkhorn normalization.
The inverse weighted constraint matrix appearing there is the same
linear-algebraic object that underlies our projection. Their work also
analyzes approximate inputs. Our additional conclusions concern exact
nonlinear changes in an ordered value group, integrality of all Taylor
coefficients, and support-controlled evaluation beyond rank one.

Spanning-tree representations of electrical projections and
transfer-current matrices belong to a well-developed theory; see
Burton and Pemantle \cite{BurtonPemantle}. We use an elementary
Cauchy--Binet and Cramer's-rule proof of the particular interpolation
identity needed here. The determinant formula, its tree interpretation,
and ordinary minimum-spanning-tree exchange arguments are not proposed
as new.

The Hahn-field and surreal background is standard; the normal-form
conventions can be compared with Gonshor \cite{Gonshor} and the
preliminaries of Berarducci and Mantova \cite{BM}. The support lemma
used below is proved from the standard word-ordering theorem of
Higman \cite{Higman}. None of these foundational ingredients is
claimed as a contribution of this article.

\subsection{Comparison with the specified repository}
The audit is pinned to the repository commit
\begin{center}\small\texttt{\pin}.\end{center}
The complete research catalogue \texttt{docs/README.md} at that pin
lists 36 reports. The nearest subjects are the Markov-generator report
and the spectral-theory report \cite{RepoCatalogue,RepoMarkov,RepoSpectral}.
The former studies normalized resolvents, positive forest formulas,
scale hierarchies, and relative stability of resolvent entries. The
latter develops finite spectral theory, valuation scales, positive
Cauchy--Binet identities, and exact field-of-definition questions.

Our shared mechanism with those reports is the absence of leading
cancellation in positive finite sums. Our object is different: a
nonlinear fixed-margin diagonal normalization, its surcomplex local
branch, and its optimal entrywise gains. A stability theorem for
$s(sI+L)^{-1}$ is not the normalization theorem proved below.
Likewise, singular-value stability does not provide the whole-polydisc
support argument or the tree-deletion sharpness statement.

The audit read the catalogue, the Markov report's detailed guide, and
the first 180 lines of the spectral report's guide. The directory
\texttt{docs/new} at the pin contains only its README, not unexamined
ZIP archives. These observations support choosing the present topic;
they are not a line-by-line audit of every article or Lean declaration.
In particular, indexed connector searches were unusable as negative
evidence: searches for \texttt{Sinkhorn} and \texttt{spanning} returned
no results, but so did the positive-control query
\texttt{MarkovResolvent}, a file name present in the repository tree.

The accompanying \texttt{SOURCE\_AUDIT.md} records the scope and search
phrases. Searches that returned irrelevant results were not treated as
evidence of absence. The appropriate priority claim is therefore:
\emph{the main formulations were not found in the sources actually
compared, and are submitted as candidate original results}. There is
no assertion that every publication or every repository line has been
excluded as a possible precedent.

\paragraph{Since the pin.}
The manuscript was placed as raw source at commit \texttt{30dfb4f}. The
statements above describe the pin and are kept as provenance. This paragraph
was added during editorial assembly at \texttt{ed88b8f}; its catalogue and
search observations describe the raw-placement snapshot. The catalogue
\texttt{docs/README.md} did not change between the pin and raw placement:
it listed 36 reports, excluding this report and the three others placed
with it. The directory
\texttt{docs/new} held only its README at the pin, as stated above. Commit
\texttt{330aa79} then added nine delivered archives, this manuscript's among
them; commit \texttt{30dfb4f} unpacked, placed and removed all nine, so that
snapshot again held only the README in this directory. The positive-control
file \nolinkurl{Surreal/Algebra/MarkovResolvent.lean} was still present. It proves
finite resolvent and projection-flag algebra for the Markov report and
contains no spanning-tree, forest-sum or Cauchy--Binet statement.
At \texttt{30dfb4f}, a full-text search of every LaTeX source in the
collection finds the words \texttt{Sinkhorn}, \texttt{bistochastic},
\texttt{Kirchhoff} and \texttt{transfer current} only in this report.
The phrase \texttt{spanning tree} occurs elsewhere only once, in the proof of
the Markov report's matrix-forest identity \cite{RepoMarkov}, where it
means a directed spanning tree rooted at a vertex. The phrase
\texttt{matrix scaling} occurs elsewhere only once, in the outline of
remaining Lean development of the Prony reconstruction report
(\nolinkurl{docs/surcomplex/prony-reconstruction-at-surreal-scales/}); that
sentence does not concern the prescribed-margin problem studied here. At that same raw-placement snapshot, no
Lean source mentioned Sinkhorn scaling, spanning trees, doubly stochastic
matrices or Kirchhoff's theorem. These searches, like the original ones,
locate words, not results stated under other names. The comparison with the
two nearest reports is made precise in \cref{scale:sec:neighbours}.

\section{Hahn conventions and the summability issue}\label{scale:sec:hahn}

\subsection{Fields, valuation, and positive leading coefficients}
Throughout the local theory, $\Gamma$ is a set-sized ordered abelian
group, written additively. It need not be divisible, Archimedean,
discrete, countable, or of finite rank. Put
\begin{equation}\label{scale:eq:fields}
 F=\R((t^\Gamma)),\qquad K=\C((t^\Gamma)).
\end{equation}
A series $x=\sum_{\gamma\in S}x_\gamma t^\gamma$ has well-ordered
support $S\subseteq\Gamma$. For $x\ne0$ define
$v(x)=\min S$ and $\lc(x)=x_{v(x)}$; set $v(0)=+\infty$.
The valuation ring and its maximal ideal are
\[
 \OO=\{x:v(x)\geq0\},\qquad \mm=\{x:v(x)>0\}.
\]
These symbols refer to $K$ unless a subscript indicates $F$.
A vector is infinitesimal when every coordinate belongs to $\mm$.
All vector dimensions in this article are ordinary finite integers.

\begin{definition}\label{scale:def:positiveleading}
A nonzero $p\in K$ is \emph{positive-leading} when $\lc(p)$ is a
positive real number. Equivalently,
$p=a t^\gamma(1+u)$ with $a\in\R_{>0}$ and $u\in\mm$.
\end{definition}
Every positive element of $F$ is positive-leading. Finite nonempty sums
and products, and inverses, preserve the positive-leading property.
For a finite positive-leading family,
\begin{equation}\label{scale:eq:nocancel}
 v\left(\sum_j p_j\right)=\min_j v(p_j).
\end{equation}
Indeed, the coefficients at the least occurring exponent add to a
positive real number. This is the only positivity mechanism required
for the local surcomplex results. It does not put an order on $K$.

\subsection{Strong summability}
\begin{definition}
A set-indexed family $(x_i)_{i\in I}$ of Hahn series is \emph{strongly
summable} if $\bigcup_i\supp(x_i)$ is well ordered and every exponent
belongs to the support of only finitely many $x_i$. Its sum is formed
coefficientwise, by finite sums at each exponent.
\end{definition}
We will never infer strong summability solely from a statement that
valuations increase strictly. Nor will a strongly summable expansion
be called a limit of finite partial sums unless the relevant
topological convergence has separately been justified.

\begin{lemma}[Positive-support word finiteness]\label{scale:lem:neumann}
If $S\subseteq\Gamma_{>0}$ is well ordered, then the additive monoid
$S^*$ generated by $S$, including $0$, is well ordered. Every
$\gamma\in\Gamma$ has only finitely many representations as a finite
nondecreasing word of elements of $S$ whose sum is $\gamma$.
Consequently it has only finitely many ordered-word representations.
\end{lemma}
\begin{proof}
Use Higman's theorem: finite words over a well-quasi-ordered alphabet
are well-quasi-ordered by subsequence embedding with componentwise
increase \cite{Higman}. A well-ordered set is such an alphabet.
If $S^*$ had a strictly descending sequence, choose a nondecreasing
word representing each of its terms. Some earlier word embeds into
a later one. Positivity of all letters and nonnegativity of the
componentwise increases imply that the earlier sum is at most the
later sum, a contradiction.

If infinitely many distinct nondecreasing words had one fixed sum,
apply the same theorem. An embedding can preserve the sum only if
there are no extra letters and no strict componentwise increases.
The two words would then be identical, again a contradiction.
Each nondecreasing word has only finitely many permutations, proving
the last assertion.
\end{proof}

\begin{lemma}[Finite-generator evaluation]\label{scale:lem:evaluation}
Let $a_1,\ldots,a_s\in\OO$ and
$R_0=\C[a_1,\ldots,a_s]\subseteq\OO$.
For every formal series
$\Phi(X)=\sum_{\nu\in\N^m}b_\nu X^\nu\in R_0[[X_1,\ldots,X_m]]$
and every $x\in\mm^m$, the family $(b_\nu x^\nu)_\nu$ is strongly
summable. Evaluation respects sums, products, and compositions with
formal series of zero constant term whenever all coefficients lie
in the same such finite-generator algebra.
\end{lemma}
\begin{proof}
Let $S$ be the union of the positive supports of the $a_j$ and the
supports of the nonzero $x_i$. It is a finite union of well-ordered
positive sets, hence well ordered. Every $a_j$ has only a possible
constant term outside $S$. Since each $b_\nu$ is a polynomial in the
$a_j$, the support of every evaluated term is contained in $S^*$.

More specifically, a contribution to $b_\nu x^\nu$ uses at least
$|\nu|$ positive letters from the supports of the $x_i$. For a fixed
exponent $\gamma$, \cref{scale:lem:neumann} bounds the length of all
words summing to $\gamma$. Thus only bounded total degrees $|\nu|$
can contribute. There are only finitely many such multi-indices.
This proves local finiteness at $\gamma$, while $S^*$ proves
well-ordering of the union of supports. If an input coordinate is zero,
terms using a positive power of it vanish and can simply be omitted.
For each of the finitely many relevant coefficients, its polynomial
expression in the $a_j$ has finitely many terms. Constant contributions
of the generators cause no extra infinite sum. Thus no uniform bound
on the polynomial degrees of all the $b_\nu$ is required.

For completeness, let $G_1,\ldots,G_\ell\in R_0[[X]]$ have zero
constant term and let $\Phi\in R_0[[Y_1,\ldots,Y_\ell]]$. Each
$G_j(x)$ lies in $\mm$: every nonzero evaluated term has positive input
degree, and the union of its supports is well ordered and positive.
At a fixed Hahn exponent, word finiteness bounds the total input degree.
Since each inner monomial has positive degree, this also bounds the outer
degree and leaves only finitely many choices of inner monomials. Expanding
their finitely many coefficients therefore licenses the rearrangement
\[
 \operatorname{ev}_x(\Phi\circ G)
   =\operatorname{ev}_{G(x)}(\Phi).
\]
The same finite-contribution argument proves compatibility with sums
and products. This proves the composition assertion, including the case
of zero inner series, without a topological limit.
\end{proof}

\begin{example}[Why integral coefficients alone do not suffice]
\label{scale:ex:badformal}
In $\Gamma=\Z^2_{\mathrm{lex}}$, let
$\delta=(0,1)$ and $\beta=(1,0)$. The formal series
\[
 \Phi(X)=\sum_{n\geq1} t^{\beta-n\delta}X^n
\]
belongs to $\OO[[X]]$, since $\beta-n\delta>0$ for every $n$.
But at $X=t^\delta$ every displayed term is $t^\beta$.
The family is not strongly summable. The coefficients do not admit
the common positive-support control in \cref{scale:lem:evaluation}.
An arbitrary-rank implicit-function proof that stops after establishing
membership in $\OO[[X]]$ therefore has a genuine gap.
\end{example}

\subsection{Restricted exponential and logarithm}
For $z\in\mm$ define
\begin{equation}\label{scale:eq:explog}
 \exp z=\sum_{n\geq0}\frac{z^n}{n!},\qquad
 \log(1+z)=\sum_{n\geq1}\frac{(-1)^{n+1}z^n}{n}.
\end{equation}
The preceding lemma applies with constant coefficient algebra $\C$.
The formal identities for these series remain valid under evaluation.
In particular they are inverse group isomorphisms
$\mm\leftrightarrow1+\mm$, and for nonzero infinitesimals
\begin{equation}\label{scale:eq:valexplog}
 v(\exp z-1)=v(z),\qquad v(\log(1+z))=v(z).
\end{equation}
Their leading terms are $z$. We also use
\begin{equation}\label{scale:eq:phi}
 \ph(z)=\frac{\exp z-1}{z}
       =\sum_{n\geq0}\frac{z^n}{(n+1)!},\qquad \ph(0)=1.
\end{equation}
Thus $\ph(z)\in1+\mm$. No exponential or logarithm at an arbitrary
infinite surcomplex argument is used anywhere.

\begin{remark}[The same support engine in the spectral report]
\label{scale:rem:support}
\Cref{scale:lem:neumann} is also Lemma~9.4 of the collection's spectral
report \cite{RepoSpectral}, proved there from Higman's theorem as well, and
\cref{scale:lem:evaluation} with no generators gives the strong-evaluation
part of Lemma~9.5 there. The generators $a_1,\ldots,a_s$ are what this
report adds, and \cref{scale:ex:badformal} shows that they cannot be replaced
by arbitrary elements of $\OO$. \Cref{scale:sec:neighbours} compares the
two reports.
\end{remark}

\section{The tree interpolation identity}\label{scale:sec:trees}

\subsection{Incidence, cuts, and weighted cycles}
Let $\GG=(V,E)$ be a finite connected graph without loops, with an
arbitrary orientation, $N=|V|\geq2$ and $m=|E|$. Choose a root
vertex. The reduced incidence matrix $B\in\Z^{(N-1)\times m}$ has
$+1$ at the initial vertex and $-1$ at the terminal vertex of each
edge, with the root row deleted. It has rank $N-1$.

For positive-leading weights $p=(p_e)$ put
\begin{equation}\label{scale:eq:projection}
 W_p=\diag(p_e),\quad L_p=BW_pB^{\mathsf T},\quad
 \Pi_p=B^{\mathsf T}L_p^{-1}BW_p,\quad H_p=I-\Pi_p.
\end{equation}
The invertibility in this formula will follow from the tree identity.
The subspace $\im B^{\mathsf T}$ is the space of potential differences,
or cuts. The kernel $\ker B$ is the space of conserved edge flows.
The kernel $\ker(BW_p)$ consists of edge functions whose weighted
flows are conserved.

For a spanning tree $T$, let $B_T$ be the square submatrix with columns
in $T$. Given $f\in K^E$, solve
$B_T^{\mathsf T}u_T=f_T$ and define
\begin{equation}\label{scale:eq:treeinterp}
 \Pi_T f=B^{\mathsf T}u_T.
\end{equation}
It extends the values of $f$ on tree edges to a potential difference
on every edge. Its value on an edge is the signed sum along the
unique tree path between that edge's endpoints. Consequently every
matrix entry of $\Pi_T$ is $0$, $1$, or $-1$.

\begin{theorem}[Weighted tree interpolation]\label{scale:thm:tree}
Write
\[
 w_T(p)=\prod_{a\in T}p_a,\qquad Z(p)=\sum_{T\in\TT}w_T(p).
\]
Then
\begin{equation}\label{scale:eq:treeaverage}
 \det L_p=Z(p)\ne0,\qquad
 \Pi_p=\sum_{T\in\TT}\frac{w_T(p)}{Z(p)}\Pi_T.
\end{equation}
Moreover
\begin{equation}\label{scale:eq:projproperties}
 \Pi_p^2=\Pi_p,\quad H_p^2=H_p,\quad
 H_pB^{\mathsf T}=0,\quad BW_pH_p=0,
\end{equation}
and both $\Pi_p$ and $H_p$ have entries in $\OO$.
\end{theorem}
\begin{proof}
Cauchy--Binet gives
\[
 \det(BW_pB^{\mathsf T})
 =\sum_{|T|=N-1}(\det B_T)^2\prod_{a\in T}p_a.
\]
An incidence minor vanishes unless $T$ is a spanning tree; for a tree
it is $\pm1$. A connected selection of $N-1$ edges is a tree;
deleting a leaf other than the root and expanding its row proves the
$\pm1$ assertion by induction. A disconnected selection has deficient
rank: a component not containing the root supplies a nonzero left-kernel
vector, its vertex indicator.
This proves the determinant formula. Every summand is positive-leading,
so \eqref{scale:eq:nocancel} proves nonvanishing.

To obtain the vector identity, set $u=L_p^{-1}BW_pf$.
For coordinate $k$, replace column $k$ of $B^{\mathsf T}$ by $f$,
calling the resulting $m\times(N-1)$ matrix $C_k$. Then $BW_pC_k$
is $L_p$ with column $k$ replaced by $BW_pf$. Cauchy--Binet and
Cramer's rule yield the following formula, where $(C_k)_T$ selects
\emph{rows} indexed by $T$, whereas $B_T$ selects columns:
\[
 u_k=\frac{1}{Z(p)}\sum_T \det(B_T)w_T(p)\det((C_k)_T)
     =\frac{1}{Z(p)}\sum_T w_T(p)(u_T)_k.
\]
In the second equality we used
$\det((C_k)_T)=\det(B_T^{\mathsf T})(u_T)_k$.
Multiplying by $B^{\mathsf T}$ proves the interpolation formula.

The projection identities follow directly from
$L_p=BW_pB^{\mathsf T}$. More explicitly, $\Pi_p$ takes values in
$\im B^{\mathsf T}$ and fixes that space, since
$\Pi_pB^{\mathsf T}=B^{\mathsf T}$. Also
$BW_p\Pi_p=BW_p$, so $\Pi_pf=0$ implies $BW_pf=0$; the converse
is immediate from its formula. Hence
\[
 \im\Pi_p=\im B^{\mathsf T},\qquad
 \ker\Pi_p=\ker(BW_p),\qquad
 \im H_p=\ker(BW_p),\qquad
 \ker H_p=\im B^{\mathsf T}.
\]
These descriptions explain the complementary projections used below.
Finally $v(w_T/Z)\geq0$ by the determinant formula and no leading
cancellation. Each $\Pi_T$ has integer entries. Finite sums therefore
show $\Pi_p,H_p\in\operatorname{Mat}_m(\OO)$.
\end{proof}

\subsection{The row formula and the diagonal coefficient}
For $e\notin T$, define the fundamental-cycle functional
\[
 \ell_{T,e}(f)=f_e-(\Pi_Tf)_e.
\]
It has coefficients $0,1,-1$, its coefficient at $f_e$ is $1$, and
it vanishes on $\im B^{\mathsf T}$. If $e\in T$, the corresponding
row of $I-\Pi_T$ is zero. Thus
\begin{equation}\label{scale:eq:rowtree}
 (H_pf)_e=\sum_{\substack{T\in\TT\\e\notin T}}
              \frac{w_T(p)}{Z(p)}\ell_{T,e}(f).
\end{equation}
In particular, with $Z_e^-(p)=\sum_{T:e\notin T}w_T(p)$,
\begin{equation}\label{scale:eq:Hdiagonal}
 (H_p)_{ee}=\frac{Z_e^-(p)}{Z(p)}.
\end{equation}
For a bridge the numerator and the entire row of $H_p$ vanish.
For other edges its valuation is exactly the gap $\kappa_e$ in
\eqref{scale:eq:overviewgap}.

The matrix $L_p^{-1}$ itself need not be integral. The finite
projection $B^{\mathsf T}L_p^{-1}BW_p$ nevertheless is integral.
This cancellation of apparent conditioning loss is the algebraic
reason to work with edge variables rather than estimate vertex
potentials using an unstructured inverse bound.

\begin{remark}[Positive tree sums elsewhere in the collection]
\label{scale:rem:trees}
That a positive-leading tree sum has the minimum tree cost as its valuation,
and the sum of the minimizers' leading amplitudes as its leading coefficient,
is the mechanism of the directed forest sums of the Markov report
\cite{RepoMarkov} and of the squared-minor sums of Gram determinants in the
spectral report
\cite{RepoSpectral}. \Cref{scale:sec:neighbours} states what is shared and
what is not.
\end{remark}

\section{Normalization on the whole infinitesimal polydisc}
\label{scale:sec:normalization}

\subsection{The nonlinear equations}
Keep the finite oriented graph and positive-leading weights from the
previous section. For $f\in\mm^E$ we seek $h\in\mm^E$ satisfying
\begin{equation}\label{scale:eq:normalizingeqs}
 h-f\in\im B^{\mathsf T},\qquad
 BW_p(\exp h-\mathbf1)=0.
\end{equation}
Exponential, multiplication, and subtraction in an edge vector are
componentwise. The corresponding normalized weights are
\begin{equation}\label{scale:eq:qdef}
 q_p(f)=p\od\exp h.
\end{equation}
The second equation says $Bq_p(f)=Bp$. On a bipartite graph oriented
from rows to columns it says exactly that every row sum and column
sum is restored. Deleting one incidence row loses no condition,
because the sum of all rows of the full incidence matrix is zero.
The first equation says that $q_p(f)$ is obtained from
$p\od\exp f$ by vertex-potential factors.

\begin{theorem}[Uniform infinitesimal normalization]
\label{scale:thm:existence}
For every $f\in\mm^E$ there is exactly one $h\in\mm^E$ satisfying
\eqref{scale:eq:normalizingeqs}. Write $h=\NN_p(f)$.
There is a unique formal series
\[
 \mathcal H_p(X)\in (X_1,\ldots,X_m)\OO[[X_1,\ldots,X_m]]^E
\]
solving the same equations formally. Its coefficients belong to the
finite-generator algebra
\begin{equation}\label{scale:eq:coeffring}
 R_p=\C[(\Pi_p)_{ab}:a,b\in E]\subseteq\OO.
\end{equation}
It evaluates strongly at every $f\in\mm^E$, and that evaluation is
$\NN_p(f)$. No lower bound on the entries $p_e$, no rank-one
hypothesis, and no divisibility of $\Gamma$ are required.
\end{theorem}
\begin{proof}
Put $R(z)=\exp z-1-z$, componentwise. Applying the complementary
projections to \eqref{scale:eq:normalizingeqs} shows that these
equations are equivalent to
\begin{equation}\label{scale:eq:fixedpoint}
 h=H_pf-\Pi_pR(h).
\end{equation}
Indeed, the cut condition is $H_ph=H_pf$, and conservation is
$\Pi_p(h+R(h))=0$. Adding these two projected identities gives
\eqref{scale:eq:fixedpoint}; the converse follows from
\cref{scale:thm:tree}. The image and kernel descriptions extend
coefficientwise to the modules over $K[[X]]$, giving the same equivalence
for the formal equations.

Now work formally in variables $X=(X_e)$. Write
$\mathcal H_p=\sum_{d\geq1}h_d$ with each $h_d$ homogeneous of degree
$d$. Since $R$ starts in degree two, the recurrence is triangular:
\begin{align}
 h_1&=H_pX,\label{scale:eq:recursionone}\\
 h_d&=-\Pi_p\sum_{k=2}^{d}\frac1{k!}
       \sum_{\substack{d_1+\cdots+d_k=d\\d_j\geq1}}
             h_{d_1}\od\cdots\od h_{d_k},\quad d\geq2.
             \label{scale:eq:recursion}
\end{align}
Every index $d_j$ on the right is less than $d$. This constructs one
and only one formal solution. Induction shows that its coefficients
are polynomials over $\Q$ in the entries of $\Pi_p$; $H_p=I-\Pi_p$
requires no extra generators. Thus they belong to $R_p\subseteq\OO$.

Apply \cref{scale:lem:evaluation} to this specific coefficient algebra.
It proves strong evaluation of the entire series at $f\in\mm^E$,
not merely formal existence. It also permits the products and
exponential substitution in \eqref{scale:eq:fixedpoint}. The evaluated
vector lies in $\mm^E$: each term has positive degree in infinitesimal
inputs, and its support lies in a well-ordered positive set. Hence
it gives a solution of \eqref{scale:eq:normalizingeqs}.

For uniqueness, let $h,k\in\mm^E$ solve \eqref{scale:eq:fixedpoint}.
Suppose $h\ne k$ and let $\nu=\min_e v(h_e-k_e)$.
There is a common $\eta>0$ not exceeding the valuations of any
nonzero coordinate of $h$ or $k$, since there are finitely many.
The formal factorization
\[
 R(z)-R(w)=(z-w)\psi(z,w),\qquad
 \psi(z,w)=\sum_{n\geq2}\frac1{n!}
     \sum_{j=0}^{n-1}z^{n-1-j}w^j\in(z,w)\C[[z,w]],
\]
together with strong evaluation gives
\[
 \min_e v(R(h_e)-R(k_e))\geq\nu+\eta.
\]
But $h-k=-\Pi_p(R(h)-R(k))$, and every entry of $\Pi_p$ is integral.
It follows that $\nu\geq\nu+\eta$, impossible in an ordered group.
This argument uses one strict valuation improvement, not convergence
of a sequence of improvements. It is therefore valid at arbitrary
rank.
\end{proof}

\subsection{Gauge factors, retraction, and local coordinates}
\begin{corollary}\label{scale:cor:gauge}
There is a unique root-normalized potential $u\in\mm^{N-1}$ with
$\NN_p(f)-f=B^{\mathsf T}u$. Thus in the matrix case the normalizing
row factors and column factors may all be chosen in $1+\mm$.
Furthermore, for $a\in\mm^{N-1}$,
\begin{equation}\label{scale:eq:gaugeinvariance}
 \NN_p(f+B^{\mathsf T}a)=\NN_p(f).
\end{equation}
The map is a retraction onto the set of $h\in\mm^E$ satisfying
$BW_p(\exp h-\mathbf1)=0$.
\end{corollary}
\begin{proof}
Existence of a potential over $K$ follows from the cut condition.
Restrict that condition to any spanning tree. The inverse of
$B_T^{\mathsf T}$ has integer entries $0,1,-1$, given by signed
root-path sums. Thus the root-normalized
potential is a finite signed sum of infinitesimal edge differences
and is itself infinitesimal. Uniqueness follows from the rank of $B$.
For an edge oriented from row $i$ to column $j$, the multiplier is
$\exp(u_i-u_j)$, with the root potential interpreted as zero.
These are row factors $\exp u_i$ and column factors $\exp(-u_j)$.
Changing $f$ by a cut does not change its class modulo cuts, so
uniqueness proves \eqref{scale:eq:gaugeinvariance}. If $h$ is already
balanced, it solves the problem with input $h$, proving the retraction
assertion.
\end{proof}

The weighted cycle coordinates are especially simple:
\begin{equation}\label{scale:eq:cyclechart}
 H_p\NN_p(f)=H_pf.
\end{equation}
Consequently $x\mapsto\NN_p(x)$ is a bijection from
$\ker(BW_p)\cap\mm^E$ onto the balanced infinitesimal set, with
inverse $h\mapsto H_ph$. Indeed $H_p\NN_p(x)=H_px=x$ on that
kernel. Conversely, for balanced $h$, gauge invariance and retraction give
$\NN_p(H_ph)=\NN_p(h)=h$; $H_ph-h$ is an infinitesimal cut.
Integrality of $H_p$ ensures that these maps preserve the infinitesimal
domains. At a balanced point $h$, the formal
linearization of conservation has kernel
$\ker(B\diag(p_e\exp h_e))$. The diagonal is invertible and $B$
has rank $N-1$, so this formal tangent space has dimension $m-N+1$
over $K$; at $h=0$ it is $\ker(BW_p)$. This is a support-controlled analytic chart in
Hahn coordinates, not an assertion about classical holomorphic
manifolds in the full surreal fine topology.

For reference, the first terms are
\begin{align}
 h_1&=H_pX,\notag\\
 h_2&=-\tfrac12\Pi_p(h_1\od h_1),\label{scale:eq:firstterms}\\
 h_3&=\tfrac12\Pi_p\bigl(h_1\od\Pi_p(h_1\od h_1)\bigr)
       -\tfrac16\Pi_p(h_1\od h_1\od h_1).\notag
\end{align}
Only the first line is the familiar implicit differentiation formula.
The theorem supplies all orders and justifies their evaluation on a
common relative infinitesimal domain.

\section{Sharp nonlinear valuation gains}\label{scale:sec:gain}

\subsection{Secant reweighting}
The decisive step turns an exact nonlinear comparison into a linear
projection with slightly altered weights.

\begin{lemma}[Exact secant identity]\label{scale:lem:secant}
Let $f,g\in\mm^E$, set
$h=\NN_p(f)$, $k=\NN_p(g)$, $r=h-k$, and $d=f-g$.
Define
\begin{equation}\label{scale:eq:secantweights}
 \widetilde p_e=p_e\exp(k_e)\ph(r_e).
\end{equation}
Then $\widetilde p_e$ has the same valuation and leading coefficient
as $p_e$, and
\begin{equation}\label{scale:eq:secantidentity}
 r=H_{\widetilde p}d.
\end{equation}
\end{lemma}
\begin{proof}
Both $\exp(k_e)$ and $\ph(r_e)$ belong to $1+\mm$, so the leading
data are unchanged. Subtract the two conservation equations:
\[
 0=B\diag(p_e\exp k_e)(\exp r-\mathbf1)
   =B\diag(\widetilde p_e)r.
\]
Also $r-d\in\im B^{\mathsf T}$. Project with
$H_{\widetilde p}$: it fixes $r$ by the first identity and kills
$r-d$ by the second. This gives \eqref{scale:eq:secantidentity}.
\end{proof}

\subsection{The optimal two-point bound}
For a tree $T$ put $c(T)=\sum_{a\in T}c_a$, with $c_a=v(p_a)$.
Define $\tau$, $\tau_e^-$, and $\kappa_e$ by
\eqref{scale:eq:overviewgap}. If $e$ is not a bridge, let
\begin{equation}\label{scale:eq:leadinggain}
 C_e=
 \frac{\displaystyle\sum_{\substack{T:e\notin T\\c(T)=\tau_e^-}}
                         \prod_{a\in T}\lc(p_a)}
      {\displaystyle\sum_{\substack{T\\c(T)=\tau}}
                         \prod_{a\in T}\lc(p_a)}\in\R_{>0}.
\end{equation}
Ties are included in both sums.

\begin{theorem}[Sharp nonlinear tree-gap law]\label{scale:thm:sharp}
For every pair $f,g\in\mm^E$ with $f\ne g$, put
$\delta=\min_a v(f_a-g_a)>0$. Then for every edge $e$,
\begin{equation}\label{scale:eq:sharpbound}
 v\bigl(\NN_p(f)_e-\NN_p(g)_e\bigr)
 =v\!\left(\frac{q_p(f)_e}{q_p(g)_e}-1\right)
 \geq\delta+\kappa_e.
\end{equation}
If $e$ is a bridge, its normalized entry is exactly $p_e$ for all
inputs. If $e$ is not a bridge and $f-g=s\mathbf e_e$ with
$0\ne s\in\mm$, then
\begin{equation}\label{scale:eq:exactsharp}
 v\bigl(\NN_p(f)_e-\NN_p(g)_e\bigr)=v(s)+\kappa_e,
\end{equation}
and, more precisely,
\begin{equation}\label{scale:eq:exactamplitude}
 \lc\!\left(
   \frac{\NN_p(f)_e-\NN_p(g)_e}{s}\right)=C_e,
 \qquad
 v\!\left(
   \frac{\NN_p(f)_e-\NN_p(g)_e}{s}\right)=\kappa_e.
\end{equation}
The same leading amplitude holds with the logarithmic difference
replaced by $q_p(f)_e/q_p(g)_e-1$.
\end{theorem}
\begin{proof}
Apply the row formula \eqref{scale:eq:rowtree} to the exact secant
identity:
\begin{equation}\label{scale:eq:secantrow}
 r_e=\sum_{T:e\notin T}
       \frac{w_T(\widetilde p)}{Z(\widetilde p)}\ell_{T,e}(d).
\end{equation}
Every cycle functional has valuation at least $\delta$, since its
coefficients are $0,1,-1$. All reweighted entries have their original
leading data, so no leading cancellation occurs in the tree sums.
In particular
\[
 v Z(\widetilde p)=\tau,\qquad
 v\bigl(w_T(\widetilde p)/Z(\widetilde p)\bigr)=c(T)-\tau.
\]
Each tree appearing in \eqref{scale:eq:secantrow} avoids $e$, hence
each summand has valuation at least $\delta+\kappa_e$.
This proves the inequality. For a bridge the sum is empty; taking
$g=0$ gives $r_e=0$ and $q_p(f)_e=p_e$.

If $d=s\mathbf e_e$, then $\ell_{T,e}(d)=s$ for every tree avoiding
$e$. The exact identity becomes
\begin{equation}\label{scale:eq:owncoordinate}
 r_e=s\,\frac{Z_e^-(\widetilde p)}{Z(\widetilde p)}.
\end{equation}
Both finite sums are nonzero and positive-leading. Their valuations
are $\tau_e^-$ and $\tau$, and their leading coefficients are the
numerator and denominator in \eqref{scale:eq:leadinggain}. This proves
\eqref{scale:eq:exactsharp} and \eqref{scale:eq:exactamplitude}, without
a genericity assumption. Finally $q_p(f)_e/q_p(g)_e=\exp r_e$, so
\eqref{scale:eq:valexplog} gives equality of valuations and preservation
of the leading coefficient for the relative output error.
\end{proof}

\begin{remark}[What ``sharp'' means]
If $\Gamma\ne\{0\}$, the quantity $\kappa_e$ is the largest gain
valid uniformly over all infinitesimal changes and all centers in this
relative polydisc: any nonzero $s\in\mm$ tests its own coordinate.
For $\Gamma=\{0\}$ there is only the zero input, so this evaluated
optimality assertion is replaced by the formal coefficient assertion below.
A particular perturbation can have a larger gain because of cycle
cancellation or because it is supported far away. The chain theorem
below gives exact examples of such stronger, location-sensitive gains.
Sharpness does not assert equality for every input direction.
\end{remark}

\subsection{All Taylor orders and a remainder certificate}
\begin{theorem}[Coefficientwise gain]\label{scale:thm:coeffgain}
For each edge $e$ and each nonzero multi-index $\nu$,
\begin{equation}\label{scale:eq:coefficientbound}
 v\bigl([X^\nu]\mathcal H_p(X)_e\bigr)\geq\kappa_e,
 \qquad
 v\bigl([X^\nu](\exp\mathcal H_p(X)_e-1)\bigr)\geq\kappa_e.
\end{equation}
For a nonbridge, the coefficient of $X_e$ has valuation exactly
$\kappa_e$ and leading coefficient $C_e$. At every center $g\in\mm^E$
the same assertions hold for the centered branch with base weights
$q_p(g)$.
\end{theorem}
\begin{proof}
Set formally $\widetilde p_e=p_e\ph(\mathcal H_p(X)_e)$.
The formal weighted Laplacian has constant term $L_p$, so its determinant
has nonzero constant term $Z(p)$ and is a unit in $K[[X]]$.
The adjugate, Cauchy--Binet and Cramer identities used for tree interpolation
therefore apply over this formal ring as well. No notion of positive-leading
$X$-series is needed. The formal secant argument gives
\[
 \mathcal H_p(X)_e=
 \sum_{T:e\notin T}
     \frac{w_T(\widetilde p)}{Z(\widetilde p)}\ell_{T,e}(X).
\]
The denominator can be written as
\[
 \frac{Z(\widetilde p)}{Z(p)}
 =\sum_T\frac{w_T(p)}{Z(p)}
          \prod_{a\in T}\ph(\mathcal H_p(X)_a).
\]
It has coefficients in $\OO$ and constant term $1$. Therefore its
formal inverse has coefficients in $\OO$. For every tree avoiding
$e$, the scalar $w_T(p)/Z(p)$ has valuation at least $\kappa_e$.
All other factors in the row formula have integral coefficients.
The first bound in \eqref{scale:eq:coefficientbound} follows. For a
bridge the row sum is empty, so the entire formal row is zero and both
bounds with $\kappa_e=+\infty$ hold. For a nonbridge,
since $\kappa_e\geq0$, every positive power of that row series also
has all its nonzero coefficients of valuation at least $\kappa_e$;
this proves the exponential bound. The coefficient of $X_e$ is
$(H_p)_{ee}=Z_e^-(p)/Z(p)$, with the asserted leading data.

For a center $g$, put $h=\NN_p(g+x)$, $k=\NN_p(g)$ and $r=h-k$.
Then $r-x\in\im B^{\mathsf T}$ and
\[
 B\diag(q_p(g))(\exp r-\mathbf1)
   =B\diag(p)(\exp h-\exp k)=0.
\]
The vector $r$ is infinitesimal, and $q_p(g)$ has the same positive leading
data as $p$. Uniqueness at this new base gives the exact recentering identity
\begin{equation}\label{scale:eq:recenter}
 \NN_p(g+x)-\NN_p(g)=\NN_{q_p(g)}(x),\qquad x\in\mm^E.
\end{equation}
Its right side is defined by its own support-controlled formal
series, so no unjustified substitution of a nonzero constant into
an arbitrary formal series is involved. The weights $q_p(g)$ have
the same valuations and leading coefficients as $p$. Apply the
proved result with that base.
\end{proof}

\begin{corollary}[Finite-degree precision certificate]
\label{scale:cor:remainder}
If $D\in\N_0$, $f\in\mm^E$ is nonzero, $\delta=\min_e v(f_e)$, and
$h^{[D]}=\sum_{d=1}^D h_d(f)$, then
\begin{equation}\label{scale:eq:remainder}
 v\bigl(\NN_p(f)_e-h^{[D]}_e\bigr)
 \geq\kappa_e+(D+1)\delta.
\end{equation}
An analogous bound holds for truncating the relative output series.
\end{corollary}
\begin{proof}
Every term of total degree $d>D$ has valuation at least
$\kappa_e+d\delta$. The family is strongly summable by the existence
theorem. Its sum cannot acquire an exponent below the common lower
bound $\kappa_e+(D+1)\delta$. The relative output series
$\exp(\mathcal H_p)-\mathbf1$ also has coefficients in $R_p$ and zero
constant term, so the same evaluation lemma and coefficient bound apply.
The case $D=0$ uses the empty truncation; bridge rows are identically zero.
\end{proof}
This is an exact certificate, but it is not a promise that finite
Taylor degree reaches every requested precision. In a higher-rank
group the set of multiples $D\delta$ can be bounded above.
The complete Hahn value still exists by its support construction.

\section{Computing the gain without summing trees}\label{scale:sec:bottleneck}

\begin{theorem}[Bottleneck replacement formula]\label{scale:thm:bottleneck}
For a nonbridge edge $e$ joining vertices $u$ and $v$, define
\[
 b_e=\min_{\pi:u\leadsto v\text{ in }\GG\setminus e}
                   \max_{a\in\pi}c_a,
\]
where it suffices to consider simple paths. Then
\begin{equation}\label{scale:eq:bottleneck}
 \kappa_e=\max(0,b_e-c_e).
\end{equation}
Thus $\kappa_e>0$ precisely when $e$ belongs to every
minimum-valuation spanning tree; $\kappa_e=0$ precisely when some
minimum tree avoids $e$.
\end{theorem}
\begin{proof}
The finite nonempty set of simple alternative paths attains its minimum.
Choose a minimum-weight spanning tree $T_0$ of $\GG\setminus e$.
The maximum edge cost on its $u$--$v$ path is $b_e$.
To prove this familiar bottleneck property, remove a maximum-cost
edge of that path. Every alternative $u$--$v$ path crosses the
resulting cut. If an alternative path had every cost below that
maximum, one of its crossing edges could replace the removed edge
and make $T_0$ cheaper, a contradiction.

Insert $e$ into $T_0$ and remove a path edge of cost $b_e$.
The resulting tree has total cost $\tau_e^-+c_e-b_e$.
On the other hand, let $U$ be any tree containing $e$. Removing
$e$ separates it into two components. A path realizing $b_e$ crosses
this cut in some edge $a$ of cost at most $b_e$. Hence
\[
 \tau_e^-\leq c(U)-c_e+c_a\leq c(U)-c_e+b_e.
\]
It follows that the least tree cost in the full graph is
\[
 \tau=\min\{\tau_e^-,\tau_e^-+c_e-b_e\}
      =\tau_e^-+\min(0,c_e-b_e).
\]
Subtracting from $\tau_e^-$ proves \eqref{scale:eq:bottleneck}.
The last assertions also follow directly from the definitions of
$\tau$ and $\tau_e^-$.
\end{proof}

All calculations in this proof use finite sums, comparisons, and
subtractions in $\Gamma$. No real embedding of the value group is
needed. A simple exact certificate for all gains consists of one
minimum spanning tree and replacement-edge data for its cuts.
An edge outside the chosen minimum tree already has gain zero.
For a tree edge, deleting it gives a cut; the cheapest other crossing
edge determines its replacement cost. The difference from the
original edge cost is its gain, or $+\infty$ when no replacement
exists. To justify this certificate for any chosen minimum tree $T$, let
$f$ be the cheapest replacement for $e\in T$. For any tree $U$ avoiding
$e$, its endpoint path crosses the cut of $T\setminus e$ in an edge $a$
with $c_a\geq c_f$. The tree $U+e-a$ and minimality of $T$ give
\[
 c(T)\leq c(U)+c_e-c_a\leq c(U)+c_e-c_f.
\]
Thus $T-e+f$ attains $\tau_e^-=c(T)-c_e+c_f$. This argument includes ties.
These are standard finite minimum-spanning-tree operations;
the new use is the precision certificate for a nonlinear Hahn
normalization.

The positive leading coefficient $C_e$ contains more information
than $\kappa_e$. Computing it requires the weighted sums of all
minimizing trees in \eqref{scale:eq:leadinggain}, or equivalent
determinant initial terms. A single minimum tree does not determine
that amplitude when there are ties.

\section{Exact nonlinear propagation through a chain}
\label{scale:sec:chain}

\subsection{The matrix and the response}
Let $n\geq2$ and choose positive infinitesimals
$a_1,\ldots,a_{n-1}\in F$, with $\gamma_i=v(a_i)>0$.
Put $a_0=a_n=0$ and
\begin{equation}\label{scale:eq:chainmatrix}
 p_{ii}=d_i=1-a_{i-1}-a_i,\qquad
 p_{i,i+1}=p_{i+1,i}=a_i,
\end{equation}
with all other entries zero. The matrix $P$ is positive on its
tridiagonal support and is bistochastic. Each $d_i$ has valuation
zero and leading coefficient $1$.

Perturb only the bottom-right kernel entry:
\[
 A_{nn}=p_{nn}\exp s,\qquad A_{ij}=p_{ij}\ \ ((i,j)\ne(n,n)),
 \qquad 0\ne s\in\mm_K.
\]
Let $Q$ be the normalization from \cref{scale:thm:existence}, and
write $\delta=v(s)$. The result below gives all nonzero response
valuations, not just a common bound.

\begin{theorem}[Additive attenuation along the chain]
\label{scale:thm:chain}
For $1\leq i\leq n-1$,
\begin{equation}\label{scale:eq:chainoffdiag}
 v\!\left(\frac{Q_{i,i+1}}{a_i}-1\right)
 =v\!\left(\frac{Q_{i+1,i}}{a_i}-1\right)
 =\delta+\sum_{k=i+1}^{n-1}\gamma_k.
\end{equation}
For the diagonal entries,
\begin{align}
 v\!\left(\frac{Q_{ii}}{d_i}-1\right)
     &=\delta+\sum_{k=i}^{n-1}\gamma_k,
            &&1\leq i<n,\label{scale:eq:chaindiag}\\
 v\!\left(\frac{Q_{nn}}{d_n}-1\right)
     &=\delta+\gamma_{n-1}.\label{scale:eq:chainlast}
\end{align}
These are exact nonlinear identities for arbitrary positive
$\gamma_i$ in the ordered group. They do not require the
$\gamma_i$ to be comparable by Archimedean size.
\end{theorem}

\subsection{Cycle coordinates and the secant matrix}
\begin{proof}
Let $r=n-1$. For $1\leq i\leq r$, define the edge vector $C_i$ to
have $+1$ at $(i,i)$ and $(i+1,i+1)$, $-1$ at $(i,i+1)$ and
$(i+1,i)$, and zero elsewhere. These are conserved flows. They are
linearly independent because each off-diagonal pair belongs to only
one $C_i$. The connected support graph has $2n$ vertices and
$3n-2$ edges, so its cycle-space dimension is $n-1$. Hence the
columns of $C=(C_1,\ldots,C_r)$ form a basis of $\ker B$.

Write the change in normalized entries as
\begin{equation}\label{scale:eq:flowcoordinates}
 j=Q-P=Cx.
\end{equation}
Then
\[
 j_{i,i+1}=j_{i+1,i}=-x_i,\qquad
 j_{ii}=x_{i-1}+x_i,
\]
where $x_0=x_n=0$.
Set $h=\NN_p(s\mathbf e_{nn})$ and
$\widetilde p_e=p_e\ph(h_e)$. The identity
$Q-P=\diag(\widetilde p_e)h$ gives
$h=\diag(\widetilde p_e)^{-1}Cx$.
Since $C^{\mathsf T}B^{\mathsf T}=0$, the cut condition implies
$C^{\mathsf T}h=C^{\mathsf T}(s\mathbf e_{nn})=s\mathbf e_r$.
Thus
\begin{equation}\label{scale:eq:chainM}
 Mx=s\mathbf e_r,\qquad
 M=C^{\mathsf T}\diag(\widetilde p_e)^{-1}C.
\end{equation}
This equation is exact even though its coefficients depend on the
unknown normalized solution. That dependence is harmless for
valuations, because the leading data of $\widetilde p$ and $p$
coincide.

The matrix $M$ is symmetric tridiagonal. Its entries are
\begin{align*}
 M_{ii}&=\widetilde p_{ii}^{-1}
        +\widetilde p_{i+1,i+1}^{-1}
        +\widetilde p_{i,i+1}^{-1}
        +\widetilde p_{i+1,i}^{-1},\\
 M_{i,i+1}&=M_{i+1,i}=\widetilde p_{i+1,i+1}^{-1}.
\end{align*}
Consequently
\begin{equation}\label{scale:eq:chainMvals}
 v(M_{ii})=-\gamma_i,\qquad v(M_{i,i+1})=0,
 \qquad M_{ii}\sim\frac{2}{a_i},\quad M_{i,i+1}\sim1.
\end{equation}
Here $u\sim v$ means $u/v\in1+\mm$. The two inverse off-diagonal
weights in $M_{ii}$ have the same positive leading coefficient and
cannot cancel.

Let $D_{a,b}$ be the determinant of the principal interval indexed
from $a$ to $b$, with the empty determinant equal to $1$.
The initial values are $D_{a,a}=M_{aa}$ and the empty determinant $1$.
For $b\geq a+1$ the continuant recurrence is
\[
 D_{a,b}=M_{bb}D_{a,b-1}-M_{b-1,b}^2D_{a,b-2}.
\]
Induction proves
\begin{equation}\label{scale:eq:continuantval}
 v(D_{a,b})=-\sum_{k=a}^b\gamma_k,
 \qquad D_{a,b}\sim\prod_{k=a}^b\frac{2}{a_k}.
\end{equation}
For an interval of length at least two, the second term in the
recurrence has valuation strictly larger than that of the first
by $\gamma_{b-1}+\gamma_b>0$. This proves both statements and also
proves that $M$ is invertible.

For $i\leq j$, cofactor expansion of a tridiagonal matrix gives
\begin{equation}\label{scale:eq:tridiaginverse}
 (M^{-1})_{ij}=(-1)^{i+j}
 \frac{\left(\prod_{k=i}^{j-1}M_{k,k+1}\right)
              D_{1,i-1}D_{j+1,r}}{D_{1,r}}.
\end{equation}
One way to see the factorization is to delete row $j$ and column
$i$ in the cofactor. Expanding successively through the intervening
rows forces the entries $M_{i,i+1},\ldots,M_{j-1,j}$; the remaining
blocks are the two displayed principal intervals. The cofactor sign
is $(-1)^{i+j}$. Symmetry gives the formula for $i>j$.
Equations \eqref{scale:eq:chainMvals}--\eqref{scale:eq:tridiaginverse}
now imply
\begin{equation}\label{scale:eq:inversechainval}
 v((M^{-1})_{ij})=\sum_{k=\min(i,j)}^{\max(i,j)}\gamma_k.
\end{equation}
In particular,
\begin{equation}\label{scale:eq:xival}
 v(x_i)=\delta+\sum_{k=i}^{r}\gamma_k,
 \qquad
 x_i\sim(-1)^{r-i}\frac{s\,a_i a_{i+1}\cdots a_r}{2^{r-i+1}}.
\end{equation}
These are exact statements about the nonlinear solution of
\eqref{scale:eq:chainM}, not substitutions into a linearization at
$s=0$.

For an off-diagonal entry, divide $-x_i$ by $a_i$ to obtain
\eqref{scale:eq:chainoffdiag}. If $1<i<n$, then
$v(x_{i-1})=v(x_i)+\gamma_{i-1}>v(x_i)$, so the two summands in
$j_{ii}=x_{i-1}+x_i$ cannot cancel at the leading exponent.
The first diagonal entry is $j_{11}=x_1$ and the last is
$j_{nn}=x_r$. Since all $d_i$ have valuation zero, the diagonal
formulas follow. For the smallest chain $n=2$, $M$ has one entry,
$x_1\sim s a_1/2$, both diagonal relative responses have valuation
$\delta+\gamma_1$, and both off-diagonal relative responses have valuation
$\delta$. This also checks the empty sums in the statement.
\end{proof}

\subsection{A rank-two example}
Take $n=3$, $a_1=t^\alpha$, $a_2=t^\beta$, with
$\alpha=(0,1)$ and $\beta=(1,0)$ in $\Z^2_{\mathrm{lex}}$.
Let $s\ne0$ be any surcomplex infinitesimal. The leading responses are
\begin{center}
\begin{tabular}{@{}lll@{}}
\toprule
Entry & Relative error, up to a factor in $1+\mm$ & Valuation\\
\midrule
$(1,1)$ & $-s\,t^{\alpha+\beta}/4$ & $v(s)+\alpha+\beta$\\
$(2,2),(3,3)$ & $s\,t^\beta/2$ & $v(s)+\beta$\\
$(1,2),(2,1)$ & $s\,t^\beta/4$ & $v(s)+\beta$\\
$(2,3),(3,2)$ & $-s/2$ & $v(s)$\\
\bottomrule
\end{tabular}
\end{center}
The perturbation reaches the adjacent small link at its original
relative scale, but its influence on the far end acquires both
link valuations. The formulas remain valid when $v(s)=\alpha$;
there is no need for integer multiples of $\alpha$ to cross $\beta$.
The far diagonal has uniform own-entry gain $\alpha$, but this remote
perturbation gains $\alpha+\beta$. The chain theorem is therefore
strictly more location-sensitive than the uniform row bound.
The signs in the table also illustrate why off-coordinate leading
responses can alternate even though the sharp own-coordinate
amplitude in \cref{scale:thm:sharp} is positive.

\section{Global positive matrix scaling and its relation to the branch}
\label{scale:sec:global}

\subsection{Classical existence and real-closed transfer}
The following background proposition explains how the local result
applies to the usual positive normalization, not only to an abstract
choice of a surcomplex branch.

\begin{proposition}[Positive scaling over real closed fields]
\label{scale:prop:global}
Let $F_0$ be a real closed ordered field. Fix a finite bipartite
support $E$ and prescribed row and column margins. Suppose a matrix
$P$ with exactly this support, strictly positive on $E$, realizes
those margins. Then every matrix $A$ strictly positive on the same
support has a unique normalized matrix
$Q=D_r A D_c$ with those margins, where $D_r,D_c$ are positive
diagonal matrices. On a connected support the factors are unique
up to multiplying $D_r$ by a common positive scalar and $D_c$ by
its inverse.
\end{proposition}
\begin{proof}
Over $\R$, consider the nonnegative matrices on $E$ with the given
margins and minimize
\[
 \sum_{e\in E} x_e\bigl(\log(x_e/A_e)-1\bigr),
 \qquad 0\log0:=0.
\]
The feasible polytope is nonempty and compact: each entry is bounded
by its row sum. The objective is strictly convex, so a minimizer
exists and is unique. It has no zero entry on $E$. Indeed, interpolate
from a putative boundary minimizer toward the strictly positive
feasible witness $P$. A zero coordinate contributes a negative
term of order $-c\,s|\log s|$, while changes in positive coordinates
are only $O(s)$. The objective would decrease for sufficiently
small positive real $s$.

At the interior minimizer, the gradient vector
$(\log(x_e/A_e))$ annihilates the tangent space of the margin
constraints. It is therefore a sum of a row potential and a column
potential. Exponentiating gives the diagonal factors. Conversely,
such a positive matrix is a constrained critical point and is the
unique minimizer. If two pairs of factors give the same matrix,
their ratios are reciprocal across each support edge; connectedness
makes them one common gauge factor. This is a standard form of
classical scaling theory \cite{Idel,SGG}.

For each fixed finite support, the existence and uniqueness assertion
can be written using only polynomial equalities, strict inequalities,
and finitely many quantified field elements: the positive witness
$P$, the input entries, the margins, the factors, and the output
entries. Completeness of the theory of real closed fields \cite{Tarski} transfers
that assertion from $\R$ to $F_0$. The real optimization proof is
not repeated inside a non-Archimedean field, and no non-Archimedean
compactness or global logarithm is assumed.
\end{proof}

In particular, the proposition applies to $F=\R((t^\Gamma))$ when
$\Gamma$ is divisible. The local results earlier require no such
hypothesis.

\subsection{An elementary multiplicative comparison}
\begin{proposition}[A finite cut bound]\label{scale:prop:cutbound}
Over any ordered field, suppose two positive kernels $A,A'$ on a
connected bipartite graph with $N$ vertices have positive diagonal
normalizations $P,P'$ to the same margins. If $\eta\geq1$ and
\[
 \eta^{-1}\leq A'_e/A_e\leq\eta\quad(e\in E),
\]
then
\begin{equation}\label{scale:eq:cutbound}
 \eta^{-N}\leq P'_e/P_e\leq\eta^N\quad(e\in E).
\end{equation}
This is a useful elementary comparison bound; optimality of the
ordinary exponent $N$ is not claimed.
\end{proposition}
\begin{proof}
Write $P=D_rAD_c$ and $P'=D'_rA'D'_c$.
Associate a positive vertex number $x_i=D'_{r,i}/D_{r,i}$ to a row
vertex and $x_j=D_{c,j}/D'_{c,j}$ to a column vertex. Then
\[
 \frac{P'_{ij}}{P_{ij}}=
       \frac{A'_{ij}}{A_{ij}}\frac{x_i}{x_j}.
\]
Sort the $N$ vertex numbers in nondecreasing order. No consecutive
ratio can exceed $\eta$. Otherwise, let $U$ be the nonempty proper
set of vertices above such a gap. On every edge from a row in $U$
to a column outside $U$, one has $P'_e-P_e>0$. On every edge
from a row outside $U$ to a column in $U$, one has $P'_e-P_e<0$.
With incidence signs, every crossing edge therefore contributes
strictly positively to the total divergence over $U$. There is a
crossing edge by connectedness. But $P$ and $P'$ have the same
margins, so that divergence is zero, a contradiction.

The ratio of the largest vertex number to the smallest is thus at
most $\eta^{N-1}$. Multiplying this bound by the kernel ratio bound
proves \eqref{scale:eq:cutbound} and its reciprocal.
\end{proof}

\begin{corollary}[The positive normalization is the local branch]
\label{scale:cor:positivebranch}
Assume $\Gamma$ is divisible and $A,A'$ are positive matrices over
$F$ on the same connected support, with fixed margins realized by a matrix
strictly positive on every edge of that support. Let $P,P'$ be their
positive normalizations.
If $A'_e/A_e\in1+\mm_F$ for every edge, put
$f_e=\log(A'_e/A_e)$. Then
\[
 P'=q_P(f).
\]
In particular \cref{scale:thm:sharp} applies to the usual positive
normalization, with the tree costs computed from the normalized
base matrix $P$.
\end{corollary}
\begin{proof}
A finite maximum of the kernel ratios, their reciprocals, and $1$
gives an $\eta\in1+\mm_F$ as in the cut bound. Its positive and
negative integer powers also lie in $1+\mm_F$. Since the infinitesimals
are order-convex, the inequalities
$\eta^{-N}\leq P'_e/P_e\leq\eta^N$ imply
$P'_e/P_e\in1+\mm_F$. The vertex numbers used in that proof may
be divided by a root value. Their largest-to-smallest ratio is bounded by
$\eta^{N-1}$, so every root-normalized ratio lies between
$\eta^{-(N-1)}$ and $\eta^{N-1}$ and thus in $1+\mm_F$.
Their restricted logarithms give an infinitesimal
potential $u$ with
\[
 \log(P'/P)-f=B^{\mathsf T}u.
\]
The margin equation is automatic. Uniqueness in
\cref{scale:thm:existence} now identifies $P'$ with $q_P(f)$.
\end{proof}
The support-feasibility hypothesis means more than positive margin numbers
and a nonnegative feasible matrix. On the support of
$\left(\begin{smallmatrix}1&1\\0&1\end{smallmatrix}\right)$, row and
column margins $(1,1)$ force the off-diagonal entry to be zero; no positive
diagonal scaling can realize them on the entire support.

The distinction between $A$ and its normalized matrix $P$ matters:
the costs in the sharp theorem are valuations of the base
\emph{normalized entries}, not, in general, valuations of the
unscaled kernel. We do not claim an equally sharp formula from
unscaled valuations alone.

\section{Extension to real linear constraints}\label{scale:sec:linear}
The nonlinear support argument and sharp gain do not fundamentally
require an incidence matrix. This gives a broader version useful
for finite exponential constraint systems.

\begin{theorem}[Basis-gap normalization]\label{scale:thm:linear}
Let $B\in\R^{r\times m}$ have full row rank and let
$p\in K^m$ have positive-leading coordinates. Consider
\[
 h-f\in\im B^{\mathsf T},\qquad
 B\diag(p_e)(\exp h-\mathbf1)=0,
 \qquad f,h\in\mm^m.
\]
There is a unique solution for every $f$, represented by a
support-controlled formal series on the whole infinitesimal
polydisc. Let $\mathcal B$ be the column bases of $B$ and define
\begin{align*}
 \tau_B&=\min_{T\in\mathcal B}\sum_{a\in T}v(p_a),\\
 \tau_{B,e}^-&=\min_{\substack{T\in\mathcal B\\e\notin T}}
                              \sum_{a\in T}v(p_a),
 \qquad \kappa_{B,e}=\tau_{B,e}^--\tau_B.
\end{align*}
An empty deletion minimum is $+\infty$. The nonlinear two-point
inequality, one-coordinate equality, and all Taylor-coefficient bounds of \cref{scale:thm:sharp,scale:thm:coeffgain}
hold with these basis gaps. A column contained in every basis has
an exactly fixed normalized coordinate. In the leading amplitude
formula, each basis term acquires the factor $(\det B_T)^2$.
\end{theorem}
\begin{proof}
For a basis $T$ put
$\Pi_Tf=B^{\mathsf T}(B_T^{\mathsf T})^{-1}f_T$.
Its coefficients are real constants and hence integral in $K$.
Cauchy--Binet and the Cramer's-rule argument give
\[
 Z_B(p)=\det(B\diag(p)B^{\mathsf T})
       =\sum_{T\in\mathcal B}(\det B_T)^2\prod_{a\in T}p_a,
\]
and
\[
 \Pi_p=\sum_{T\in\mathcal B}
 \frac{(\det B_T)^2\prod_{a\in T}p_a}{Z_B(p)}\Pi_T.
\]
Every determinant square in this sum is a strictly positive real
constant. The denominator is positive-leading, the scalar ratios
are integral, and $\Pi_p,H_p$ have integral entries. The recurrence
and finite-generator proof in \cref{scale:thm:existence} therefore
apply unchanged.

For a coordinate $e\in T$, the $e$th row of $I-\Pi_T$ is zero.
For $e\notin T$, the $e$th row is a real linear functional with
coefficient $1$ at $f_e$. Its value on an infinitesimal vector of
minimum valuation $\delta$ has valuation at least $\delta$.
The secant identity is purely linear algebra plus the exponential
identity and remains valid. Repeating its rowwise proof gives the
basis-gap bound and exact equality for $f-g=s\mathbf e_e$.
The same formal denominator argument proves the coefficient bounds.
The leading determinant-square factors are exactly those displayed
in $Z_B$. If $e$ belongs to every basis (a coloop), every row in the
formula for $(H_pf)_e$ vanishes, so its normalized coordinate is exactly
$p_e$. A zero column belongs to no basis and has gap zero; its row in $\Pi_p$
is zero and its logarithmic response is precisely $f_e$.
The convention also covers $r=0$: the unique basis is empty, its determinant
and weight are $1$, and the normalization is the identity. If $m=0$,
full row rank forces $r=0$ and there are no coordinate assertions.
\end{proof}
The hypothesis that $B$ has real constant entries is explicit.
This theorem does not assert the same basis-cost formula for an
arbitrary Hahn-valued constraint matrix, where minors and basis
interpolation coefficients can carry additional valuations. For example,
take $B=(1,\ t^\beta)$, $\beta>0$, and $p=(1,1)$. Both column bases
have entry-cost sum zero, but
\[
 (H_p)_{11}=\frac{t^{2\beta}}{1+t^{2\beta}}
\]
has valuation $2\beta$, rather than the zero predicted by the formula
that ignores the constraint minors.

\section{Surreal specialization and preservation of the support group}
\label{scale:sec:surreal}

\subsection{No ramification in the relative infinitesimal problem}
\begin{corollary}[Exact local support-group descent]
\label{scale:cor:descent}
Let $\Lambda\leq\Gamma$ be the subgroup generated by the supports
of all base weights $p_e$ and all perturbations $f_e$.
Then the normalized weights, logarithmic responses, and
root-normalized infinitesimal vertex potentials belong to
$\C((t^\Lambda))$. No divisible hull, fractional exponent, or
larger ordered group is needed for the local normalization.
\end{corollary}
\begin{proof}
The Hahn field $\C((t^\Lambda))$ contains all the input entries
and is closed under the finite field operations defining $\Pi_p$.
The finite-generator evaluation argument can be carried out inside
that same field. It yields $h$ there; restricted exponentiation and
multiplication give $q$ there. Solving for the potential on a
spanning tree uses only a constant integral inverse matrix.
\end{proof}
This is a local statement. Global scaling over a nondivisible
support group can require ramification. For example, over
$\R((t^\Z))$ the positive matrix
\begin{equation}\label{scale:eq:ramificationexample}
 A=\begin{pmatrix}1&t\\1&1\end{pmatrix}
\end{equation}
cannot be scaled within that field to a bistochastic matrix.
Such a matrix would have the form
$Q=\left(\begin{smallmatrix}x&1-x\\1-x&x\end{smallmatrix}\right)$,
and invariance of the cross-ratio would give
$x^2/(1-x)^2=t^{-1}$. Thus
$2v(x)-2v(1-x)=-1$, impossible for integer valuations.
In the half-integer extension the solution is
$x=(1+t^{1/2})^{-1}$.
The contrast with \cref{scale:cor:descent} is exact: normalization
near an already normalized positive-leading base is unramified,
even though arbitrary global normalization need not be.

\subsection{Transport to surreal and surcomplex numbers}
Conway normal form identifies a surreal number with a real series
in monomials $\omega^x$ with set-sized reverse-well-ordered support;
see \cite{Gonshor,BM}. For a set-sized ordered subgroup
$\Lambda\leq(\No,+)$, the map
\begin{equation}\label{scale:eq:surrealembedding}
 \sum_{\gamma}a_\gamma t^\gamma
 \longmapsto\sum_\gamma a_\gamma\omega^{-\gamma}
\end{equation}
is an ordered field embedding on real coefficients and a field
embedding on complex coefficients into $\No[i]$. It preserves
normal-form strong sums and the restricted exponential and logarithm.
The sign in $\omega^{-\gamma}$ is important: increasing valuation
must correspond to decreasing monomial size.

Given finitely many surreal or surcomplex input entries, take the
subgroup generated by the negatives of their normal-form exponents.
It is a set. Their entire supports lie in that one workspace, and
\cref{scale:cor:descent} keeps the answer there. Therefore all local
theorems are genuine theorems about finite surreal and surcomplex
matrices, not statements about formal symbols awaiting an
interpretation.

For global positive scaling, include a feasible witness positive on every
edge of the prescribed support and pass to the divisible hull of the input support group. The resulting
real Hahn field is real closed, so \cref{scale:prop:global} gives
existence there and hence in $\No$. Uniqueness in $\No$ follows by
placing any two proposed finite answers in one larger set-sized
real-closed workspace. This avoids treating a proper class as an
ordinary set-theoretic model.

Order-preserving group embeddings commute with the local construction,
because they preserve supports, field operations, and strongly
summable evaluation. Complex conjugation also commutes with it:
\[
 \overline{\NN_p(f)}=\NN_{\overline p}(\overline f).
\]
In particular, a real base and real perturbation produce a real
answer. None of these statements concerns the Berarducci--Mantova
scalar derivation; the derivatives implicit in Taylor coefficients
are formal derivatives with respect to the edge parameters.

\section{Examples, boundaries, and failure modes}\label{scale:sec:boundaries}

\subsection{A closed-form square example}
Let $b=t^\gamma$ with $\gamma>0$, $a=1-b$, and
\[
 P=\begin{pmatrix}a&b\\b&a\end{pmatrix}.
\]
Multiply only the $(1,1)$ kernel entry by $1+z$, where
$0\ne z\in\mm_K$. The infinitesimal square-root branch
$r=(1+z)^{1/2}\in1+\mm$ is defined by its binomial series. The
normalized answer is
\begin{equation}\label{scale:eq:squareexact}
 Q=\begin{pmatrix}x&1-x\\1-x&x\end{pmatrix},\qquad
 x=\frac{ar}{b+ar}.
\end{equation}
The row and column sums are $1$, and
$x^2/(1-x)^2=(a/b)^2(1+z)$ is the kernel cross-ratio. All entries
are relatively infinitesimally close to their base values, so
uniqueness identifies this formula with the constructed branch.
The relative errors simplify exactly to
\begin{equation}\label{scale:eq:squareerrors}
 \frac{x}{a}-1=\frac{b(r-1)}{b+ar},\qquad
 \frac{1-x}{b}-1=\frac{a(1-r)}{b+ar}.
\end{equation}
Thus diagonal entries gain $\gamma$ in valuation while the two
small off-diagonal entries gain zero. The square graph has the
same gaps: deleting a diagonal forces a tree of cost $2\gamma$
instead of $\gamma$, while a minimum tree can avoid either
off-diagonal edge. Notice that small normalized entries are not
necessarily unstable in \emph{relative} valuation.

\subsection{Positivity cannot simply be removed}
\begin{example}[Large formal derivative after phase cancellation]
\label{scale:ex:phasejacobian}
On a square support, take three weights equal to $1$ and the fourth
$w=-1/3+t^\beta$, with $\beta>0$. Then
\[
 Z=1+3w=3t^\beta,\qquad (H_p)_{44}=\frac{1}{3t^\beta}.
\]
All four entry valuations are zero, but a formal derivative
coefficient has valuation $-\beta$. The positive-leading assumption
fails at $w$, and cancellation in the tree denominator is exactly
the obstruction. This disproves the naive assertion that the same
integral-coefficient theorem follows from nonvanishing entries and
their valuations alone.
\end{example}

\begin{example}[Two distinct infinitesimal surcomplex solutions]
\label{scale:ex:phasebranches}
The failure can be nonlinear, not just a poor derivative bound.
Take
\[
 P=\begin{pmatrix}1&1\\1&-1/3\end{pmatrix},\qquad
 A=\begin{pmatrix}1&1\\1&-(1/3)\exp s\end{pmatrix},
 \qquad s=t^{2\delta},\quad\delta>0.
\]
Every matrix on this support with the same margins as $P$ is of the
form
\[
 Q(z)=\begin{pmatrix}1+z&1-z\\1-z&-1/3+z\end{pmatrix}.
\]
The cross-ratio equation for diagonal equivalence with $A$ is
\begin{equation}\label{scale:eq:phasequadratic}
 (1+\tfrac13\exp s)z^2
 +\tfrac23(1-\exp s)z
 +\tfrac13(\exp s-1)=0.
\end{equation}
Its discriminant is exactly $\tfrac{16}{9}(1-\exp s)$ and hence has
leading term $-(16/9)t^{2\delta}$.
It therefore has two square roots in $K$, obtained from
$\pm(4i/3)t^\delta$ times a binomial series in $1+\mm$.
The two roots of \eqref{scale:eq:phasequadratic} satisfy
\[
 z_+\sim\tfrac{i}{2}t^\delta,
 \qquad z_-\sim-\tfrac{i}{2}t^\delta.
\]
Both yield matrices relatively infinitesimally close to $P$, and
both are diagonal-equivalent to $A$. For a full $2\times2$ support,
equality of the nonzero cross-ratios is sufficient: solve three
row/column factor equations, and the fourth follows. Here the
factors can all be chosen in $1+\mm$ because all entry ratios are
in $1+\mm$.
Thus there are two infinitesimal solutions, and an input change of
valuation $2\delta$ produces changes of valuation $\delta$.
The uniqueness and no-loss conclusions genuinely fail without
an appropriate positivity or noncancellation hypothesis. Positive
leading coefficients are a sufficient condition, not claimed to
be a necessary one.
\end{example}

\subsection{Fixed margins are essential}
Consider the tree support of
\[
 P=\begin{pmatrix}1&t^\gamma\\0&1\end{pmatrix},\qquad\gamma>0.
\]
Its row margins are $(1+t^\gamma,1)$ and its column margins are
$(1,1+t^\gamma)$. Keep the columns fixed, but change the rows to
$(1+t^\gamma+s,1-s)$ with $s=t^\delta$ and $0<\delta<\gamma$.
The unique matrix on this support with the new margins is
\[
 P'=\begin{pmatrix}1&t^\gamma+s\\0&1-s\end{pmatrix}.
\]
All margin changes are relative infinitesimals, but
\[
 v(P'_{12}/P_{12}-1)=\delta-\gamma<0.
\]
Every edge of this graph is a bridge. The exact bridge rigidity
of \cref{scale:thm:sharp} holds for \emph{fixed} margins, not for
this different problem.

\subsection{No topological convergence claim is hidden in the proof}
In $\Z^2_{\mathrm{lex}}$, with $\delta=(0,1)$ and $\beta=(1,0)$,
the strongly defined exponential
$\exp(t^\delta)=\sum_{n\geq0}t^{n\delta}/n!$ is not the limit of its
finite Taylor sums in the intrinsic valuation topology. The error
after degree $D$ has valuation $(D+1)\delta<\beta$ for every $D$.
The error never enters the neighborhood $\{x:v(x)>\beta\}$.
This does not affect existence of the Hahn sum, which is a
coefficientwise support statement.

Accordingly, we have proved neither convergence of the alternating
Sinkhorn iteration in an arbitrary-rank valuation topology nor
convergence of set-indexed sequences in the full surreal fine
topology. Our normalization is specified by exact algebraic and
strong-series identities. Algorithmic convergence and exact
normalization are different assertions.

\section{Verification, implementation, and remaining questions}
\label{scale:sec:verification}

\subsection{What the supplied program checks}
The package contains \texttt{code/verify.py}, using Python and SymPy
with exact integers, rational functions, and rational coefficients.
It does not use floating-point evidence for a valuation identity.
The recorded run used Python~3.13.5, SymPy~1.14.0, and random seed
20260922. It passed the following finite checks:
\begin{center}
\begin{tabular}{@{}L{0.39\textwidth}L{0.53\textwidth}@{}}
\toprule
Check & Recorded scope\\
\midrule
Tree interpolation and projections &
60 connected bipartite graphs; 348 spanning trees in those graphs;
2,056 projection matrix entries.\\[3pt]
Deletion-gap versus minimax path &
458 graph/cost instances from 229 connected support graphs;
2,612 edge-gap comparisons, including 1,266 bridges;
integer and lexicographic integer-pair costs.\\[3pt]
Formal normalization recurrence &
Two exact constant-weight examples through degree 7, and two
Laurent-weight square examples through degree 6;
conservation, cut identities, and coefficient gain bounds.\\[3pt]
Chain inverse valuations &
Chain sizes $n=2,3,4,5,6$, checking 55 inverse entries.\\[3pt]
Closed-form square &
Three exact rational identities for the cross-ratio and relative
errors.\\
\bottomrule
\end{tabular}
\end{center}
The full machine-readable result is
\texttt{data/verification.json}; the console record is also included.
The graph family consists of all connected spanning subgraphs of
$K_{2,2}$, $K_{2,3}$, and $K_{3,3}$, with a fixed-seed sample of 60
used for the more expensive matrix identities. All 229 graphs are
used for the two cost-group gap comparisons. Counts include repeated
underlying graphs with different cost groups where explicitly stated.

The program writes \texttt{data/verification.json} next to its own
\texttt{code/} directory, so a run replaces the recorded file; it should be
run on a copy of this directory. During editorial assembly, such a run under
Python~3.14.4 and SymPy~1.14.0 passed in about twenty seconds and
reproduced every count above. The file it wrote differed from the recorded
one only in the Python version field and in line endings. The main-text
proof review reran the unchanged suite on a scratch copy under Python~3.13.14
and SymPy~1.14.0; it passed and matched the historical JSON except for the
Python version. The historical files were not overwritten.

These checks are independent corroboration of finite formulas, not
a proof of strong Hahn summability, arbitrary-rank existence, or
surreal transport. They do not constitute a Lean verification of any
new theorem. The general proofs are the arguments in the preceding
sections; their imported inputs are identified in the bibliography.

\subsection{A reproducible symbolic workflow}
Given exact base data $p$ and a perturbation $f$, the following finite
and formal tasks are separated.
First form the reduced incidence matrix and the rational projection
$\Pi_p$. Compute the costs $v(p_e)$ and tree-deletion gains to obtain
a precision certificate. Generate the homogeneous terms from
\eqref{scale:eq:recursionone}--\eqref{scale:eq:recursion} only to the
required total degree. The evaluation lemma supplies the full
mathematical denotation; \cref{scale:cor:remainder} supplies the
precision of the finite output.

This is an exact-operation specification, not a universal computable
representation of all surreal numbers. Arbitrary Hahn coefficients
need not come with an effective leading exponent, equality test, or
support oracle. The program uses finite graphs and rational-function
examples where these questions are decidable by explicit algebra.
At higher rank, finite Taylor degree may never reach a requested
valuation threshold, as already illustrated. A practical
implementation must distinguish a formal exact expression, a
support-truncated representation, and a finite Taylor approximation.

A Lean implementation could be modular: prove the incidence-minor and
Cauchy--Binet identities over fields; isolate the positive-leading
no-cancellation lemma; formalize the triangular coefficient recursion;
then connect it to a proved strong-summability evaluation theorem.
No assertion is made that these proposed new modules are currently
implemented in the repository.

\subsection{Questions not settled here}
The principal theorems answer the concrete local stability question
posed in \cref{scale:sec:intro}. They leave other questions open.
For a general support graph and two \emph{different} edges, a complete
sharp formula for the propagation gain from one prescribed input
edge to a prescribed output edge would need to control cancellations
between signed tree-path contributions. The chain theorem solves
one structured family, not this entire problem.

A second problem is to classify arbitrary surcomplex base weights
for which the whole-polydisc branch remains unique and integral.
Positive leading coefficients are sufficient, and the examples above
show that they cannot simply be omitted. They are not an exact
classification of all possible phases or all projection matrices.

Finally, this article does not determine the convergence behavior or
iteration complexity of alternating normalization over every
higher-rank valued field. Nor does it supply a global surcomplex
matrix-scaling theorem for arbitrary complex kernels. These are
separate targets; neither is hidden inside the word ``analytic''
used for our support-controlled local branch.

\section{Conclusion}\label{scale:sec:conclusion}
Three results form the core of this article. First, a positive-leading
Hahn matrix has a unique surcomplex normalization on the entire
relative infinitesimal polydisc, with a finite-generator support proof
that works at arbitrary rank. Second, a finite spanning-tree deletion
gap is a uniform nonlinear gain in each entry, optimal when
$\Gamma\ne\{0\}$, with matching Taylor-coefficient bounds and
one-coordinate leading amplitudes.
Third, a chain admits explicit nonlinear propagation laws in which
successive link valuations add.

The mathematical mechanism is not a general promise that replacing
real scalars by surreal scalars preserves analysis. It is a specific
combination of finite positive determinant sums, weighted projections,
a secant identity, and support-controlled formal algebra. This
combination both gives the theorems and states precisely where they
can fail. The local answer remains in the exact input support group;
the passage to surreal and surcomplex numbers changes the available
scales without changing the proof.

\appendix
\section{Proof dependency map}\label{scale:app:dependencies}
\begin{center}
\begin{tabular}{@{}L{0.31\textwidth}L{0.61\textwidth}@{}}
\toprule
Statement & Dependencies\\
\midrule
Evaluation lemma &
Higman's word theorem; positive ordered-group sums;
finite-generator coefficient algebra.\\[3pt]
Tree interpolation &
Incidence minors; Cauchy--Binet; Cramer's rule;
positive-leading finite sums.\\[3pt]
Whole-polydisc existence &
Integral tree projection; triangular formal recursion;
evaluation lemma; one-step strict valuation uniqueness.\\[3pt]
Sharp nonlinear gain &
Existence; exact secant reweighting; rowwise tree formula;
no leading cancellation in both tree sums.\\[3pt]
Taylor coefficient gain &
Formal secant identity; integral inverse of a formal unit;
finite tree-weight valuation bounds.\\[3pt]
Chain propagation &
Existence; cycle basis; exact secant weights; continuant
and cofactor identities; strict valuation dominance.\\[3pt]
Positive global interpretation &
Classical real scaling; real-closed transfer; finite cut bound;
local uniqueness.\\[3pt]
Surreal specialization &
Conway normal form and its Hahn arithmetic;
set-sized support-group localization.\\
\bottomrule
\end{tabular}
\end{center}
Only the global-positive interpretation uses real closedness. Only
the last step uses a special property of surreal numbers beyond
being realizable as Hahn normal forms. None of the main local
proofs uses spherical completeness, Banach contraction, a global
surcomplex exponential, or a surreal derivation.

\section{Build and artifact record}\label{scale:app:build}
From this directory, build the report with
\begin{verbatim}
latexmk -pdf -halt-on-error -interaction=nonstopmode article.tex
\end{verbatim}
and run the checks on a copy of the directory (\cref{scale:sec:verification}),
from the root of that copy, with
\begin{verbatim}
python -m pip install -r data/requirements.txt
python code/verify.py
\end{verbatim}
The bibliography is internal; no external database or figure file
is required. Compilation verifies typesetting,
not mathematical validity. The delivered PDF was rendered and
visually inspected separately from the exact algebra checks.

As delivered, the manuscript was built as
\texttt{surreal\_matrix\_scaling.tex}, and its \texttt{build.sh}, then at the
top of the package, changed to its own directory before running the program
and the build. That script is kept unchanged in \texttt{code/}, where it
changes to \texttt{code/} and looks for \texttt{code/verify.py} and
\texttt{surreal\_matrix\_scaling.tex} there, neither of which exists. It
therefore does not work as placed; use the commands above.

The article, source audit, code, and recorded test outputs belong to
this one contribution. The package contained no copies of the
repository's existing reports, and its delivery changed no repository
file; \cref{scale:sec:provenance} records how it was placed.

\begin{thebibliography}{99}
\small
\raggedright
\bibitem{BM}
A. Berarducci and V. Mantova,
\emph{Surreal numbers, derivations and transseries},
Journal of the European Mathematical Society \textbf{20} (2018),
339--390. \url{https://doi.org/10.4171/JEMS/769}.
Preprint: \url{https://arxiv.org/abs/1503.00315}.

\bibitem{BurtonPemantle}
R. Burton and R. Pemantle,
\emph{Local characteristics, entropy and limit theorems for spanning
trees and domino tilings via transfer-impedances},
Annals of Probability \textbf{21} (1993), 1329--1371.
\url{https://arxiv.org/abs/math/0404048}.

\bibitem{Eisenberger}
M. Eisenberger, A. Toker, L. Leal-Taix\'e, F. Bernard and D. Cremers,
\emph{A unified framework for implicit Sinkhorn differentiation},
CVPR 2022; arXiv:2205.06688 (2022).
\url{https://arxiv.org/abs/2205.06688}.

\bibitem{Gonshor}
H. Gonshor, \emph{An Introduction to the Theory of Surreal Numbers},
London Mathematical Society Lecture Note Series 121,
Cambridge University Press, 1986.

\bibitem{Higman}
G. Higman, \emph{Ordering by divisibility in abstract algebras},
Proceedings of the London Mathematical Society, third series,
\textbf{2} (1952), 326--336.
\url{https://doi.org/10.1112/plms/s3-2.1.326}.

\bibitem{Idel}
M. Idel, \emph{A review of matrix scaling and Sinkhorn's normal form
for matrices and positive maps}, arXiv:1609.06349v1 (2016).
\url{https://arxiv.org/abs/1609.06349v1}.

\bibitem{SGG}
M. Sharify, S. Gaubert and L. Grigori,
\emph{Solution of the optimal assignment problem by diagonal scaling
algorithms}, arXiv:1104.3830, version 2 (2013; first version 2011).
\url{https://arxiv.org/abs/1104.3830v2}.

\bibitem{Tarski}
A. Tarski, \emph{A Decision Method for Elementary Algebra and Geometry},
prepared with the assistance of J. C. C. McKinsey,
RAND report R-109, 1951.
\url{https://www.rand.org/pubs/reports/R109.html}.

\bibitem{RepoCatalogue}
V. Reshetnikov, \emph{Surreal repository: the surreal and surcomplex
reports}, research catalogue \texttt{docs/README.md},
commit \texttt{\shortpin}, inspected 22 September 2026.
\repo. Full commit identifier is recorded in
\cref{scale:sec:comparison} and in \texttt{SOURCE\_AUDIT.md}.

\bibitem{RepoMarkov}
\emph{Surreal Markov generators at every valuation scale},
AI-assisted research report in the Surreal repository,
\texttt{docs/surreal/markov-generators-at-every-scale/}.
Detailed README inspected at commit \texttt{\shortpin}; the article was
compared during editorial assembly at commit \texttt{ed88b8f}; cited
statements rechecked at \texttt{034ab96} during the proof review.
\repo.

\bibitem{RepoSpectral}
\emph{Surcomplex spectral theory: singular values, infinitesimal rank,
multiscale stability, exact determinantal ramification, and Hankel
square-class profiles}, AI-assisted research report in the Surreal
repository, \texttt{docs/surcomplex/spectral-theory/}.
README lines 1--180 inspected at commit \texttt{\shortpin}; the article
was compared during editorial assembly at commit \texttt{ed88b8f}; cited
statements rechecked at \texttt{034ab96} during the proof review.
\repo.
\end{thebibliography}
\end{document}
